\documentclass{amsart}
\usepackage{amsmath, amsthm, amsfonts, amssymb, color,enumerate}
\usepackage[colorlinks, citecolor=blue,pagebackref,hypertexnames=false]{hyperref}


\begin{document}

\newtheorem{theorem}{Theorem}[section]
\newtheorem{proposition}[theorem]{Proposition}
\newtheorem{coro}[theorem]{Corollary}
\newtheorem{lemma}[theorem]{Lemma}
\newtheorem{definition}[theorem]{Definition}
\newtheorem{nota}[theorem]{Notation}
\newtheorem{assum}{Assumption}[section]
\newtheorem{example}[theorem]{Example}
\newtheorem{remark}[theorem]{Remark}


\title[Schatten class properties of Dunkl--Riesz transform commutators]{Weyl--Chamber geometry, Poincar\'e inequality meets Schatten class properties of Dunkl--Riesz transform commutators}



\author{Zhijie Fan}
\address{Zhijie Fan, School of Mathematical Sciences, South China Normal University, Guangzhou 510631, China}
\email{fanzhj3@mail2.sysu.edu.cn}


\author{Ji Li}
\address{Ji Li, Department of Mathematics, Macquarie University,NSW2109, Australia}
\email{ji.li@mq.edu.au}



\author{Dmitriy Zanin}
\address{School of Mathematics and Statistics, HNP-LAMA, Central South University, Changsha 410075, China}
\email{d.zanin@csu.edu.cn}






\date{}

 \subjclass[2010]{47B10, 42B20, 43A85}
\keywords{Schatten class, Riesz transform commutators, Dunkl operator, Besov space, Nearly weakly orthogonal sequence}




\begin{abstract}

\end{abstract}

\maketitle


\tableofcontents


\section{Introduction}

\begin{theorem}\label{main theorem} Let $R$ be a normalised root system in $\mathbb{R}^N,$ $N\geq 2,$ and let $G$ be the corresponding finite reflection group. Let $k:R\to\mathbb{R}_+$ be $G$-invariant multiplicity function. Let $R_{{\rm Dunkl},j}$ be the $j$-th Dunkl-Riesz transform corresponding to $R$ and $k.$ Let $w$ be the Dunkl weight corresponding to $R$ and $k.$ Let $f\in L_1^{{\rm loc}}(\mathbb{R}^N).$
\begin{enumerate}[{\rm (i)}]
\item\label{mta} If $f\in\dot{W}^{\frac{N}{p},p}(\mathbb{R}^N),$ $p>N,$ then $[R_{{\rm Dunkl},j},M_f]\in\mathcal{L}_p(L_2(\mathbb{R}^N,w(x)dx)).$ Furthermore,
$$\|[R_{{\rm Dunkl},j},M_f]\|_{\mathcal{L}_p(L_2(\mathbb{R}^N,w(x)dx))}\leq c_{p,R,k}\|f\|_{\dot{W}^{\frac{N}{p},p}(\mathbb{R}^N)};$$
\item\label{mtb} If $[R_{{\rm Dunkl},j},M_f]\in\mathcal{L}_p(L_2(\mathbb{R}^N,w(x)dx)),$ $p>N,$ then $f\in\dot{W}^{\frac{N}{p},p}(\mathbb{R}^N).$ Furthermore,
$$\|[R_{{\rm Dunkl},j},M_f]\|_{\mathcal{L}_p(L_2(\mathbb{R}^N,w(x)dx))}\leq c_{p,R,k}^{-1}\|f\|_{\dot{W}^{\frac{N}{p},p}(\mathbb{R}^N)};$$
\item\label{mtc} If $f\in\dot{W}^{1,N}(\mathbb{R}^N),$ then $[R_{{\rm Dunkl},j},M_f]\in\mathcal{L}_{N,\infty}(L_2(\mathbb{R}^N,w(x)dx)).$ Furthermore,
$$\|[R_{{\rm Dunkl},j},M_f]\|_{\mathcal{L}_{N,\infty}(L_2(\mathbb{R}^N,w(x)dx))}\leq c_{R,k}\|f\|_{\dot{W}^{1,N}(\mathbb{R}^N)};$$
\item\label{mtd} If $[R_{{\rm Dunkl},j},M_f]\in\mathcal{L}_{N,\infty}(L_2(\mathbb{R}^N,w(x)dx)),$ then $f\in\dot{W}^{1,N}(\mathbb{R}^N).$ Furthermore,
$$\|[R_{{\rm Dunkl},j},M_f]\|_{\mathcal{L}_{N,\infty}(L_2(\mathbb{R}^N,w(x)dx))}\geq c_{R,k}^{-1}\|f\|_{\dot{W}^{1,N}(\mathbb{R}^N)};$$
\end{enumerate}
\end{theorem}

\section{Preliminaries}
\setcounter{equation}{0}
\subsection{Harmonic analysis in the Dunkl setting}
Let $\mathbb{R}^N$ be the Euclidean space equipped with the usual scalar product $\langle \cdot,\,\cdot\rangle $ and the induced norm $\|\cdot\|$. For a nonzero vector $v\in \mathbb{R}^N$, the reflection $\sigma_v$ with respect to the hyperplane $v^{\perp}$ orthogonal to $v$ is defined by
\begin{align*}
	\sigma_v x:= x-2\frac{\langle x,v\rangle}{\|v\|^2}v.
\end{align*}
A finite set $R\subset \mathbb{R}^N\backslash \{0\}$ is called a root system if $\sigma_v(R)=R$ for every $v\in R$. A root system $R$ is called a normalized reduced root system if  $\|v\|^2=2$ for every $v\in R$.
The finite group $G$ generated by the reflections $\{\sigma_v\}_{v\in R}$ is called the Weyl group (reflection group) of the root system $R$.
Let $k: R \to [0,\infty)$ be a multiplicity function, which is a $G$-invariant function  fixed throughout the paper. 


Define $d(x,y)=\min_{g\in G}|x-gy|.$ Recall from \cite[Chapter VII, proof of Theorem 2.12]{MR514561} that $d(x,y)=|x-y|$ iff $x$ and $y$ belong to the same chamber.

Define the Dunkl measure $dw$ on $\mathbb{R}^N$ by
$$
dw(x) :=\prod_{v \in R} |\langle x,v\rangle|^{k(v)} dx,
$$
where $dx$ stands for the Lebesgue measure on $\mathbb{R}^N$. Let $B(x,r):=\{y\in \mathbb{R}^N:\|x-y\|< r\}$ be the Euclidean ball centred at $x$ with radius $r>0$. Let $\mathbf{N}:=N+\sum_{v\in R}k(v)$. Note that
$$w(B(tx,tr))=t^{\mathbf{N}}w(B(x,r)),\ {\rm for}\ x\in\mathbb{R}^N,\ t,r>0.$$
Furthermore,
\begin{align}\label{twosideinequ}
  w(B(x,r))\sim r^N \prod_{v \in R}(|\langle x, v \rangle |+r)^{k(v)},
\end{align}
where the implicit constant only depends on $N,R$ and $k$.  It follows from \eqref{twosideinequ} that there exists a constant $C\geq 1$ such that for every $x\in \mathbb{R}^N$,
\begin{align}\label{doublinginequality}
C^{-1}\left(\frac{r_2}{r_1}\right)^N\leq \frac{w(B(x,r_2))}{w(B(x,r_1))}\leq C\left(\frac{r_2}{r_1}\right)^{\mathbf{N}},\ {\rm for}\ 0<r_1<r_2,
\end{align}
so the numbers $\mathbf{N}$ and $N$ are called the upper and lower homogeneous dimensions, respectively.

The Dunkl operators $T_\xi$ are defined by
$$T_\xi f(x):=\partial_\xi f(x)+\sum_{v\in R}\frac{k(v)}{2}\langle v,\xi\rangle\frac{f(x)-f(\sigma_v(x))}{\langle v,x\rangle},$$
where $\partial_\xi f$ denotes the directional derivative of $f$ in the direction $\xi$. Furthermore, the Dunkl Laplacian is defined by
$$\Delta_{{\rm Dunkl}}f(x):=\sum_{j=1}^N T_{e_j}^2f(x)=\Delta_{\mathbb{R}^N}f(x)+\sum_{v\in R}k(v)\delta_v f(x),$$
where $\Delta_{\mathbb{R}^N}$ is the Laplacian on $\mathbb{R}^N$ and $\delta_v$ is defined by
$$\delta_v f(x):=\frac{\partial_v f(x)}{\langle v,x\rangle}-\frac{f(x)-f(\sigma_v(x))}{\langle v,x\rangle^2}.$$
In this article, we consider the Dunkl-Riesz transforms defined by $R_{{\rm Dunkl},j}:=T_{e_j}(-\Delta_{{\rm Dunkl}})^{-1/2}$, $j=1,2,\ldots,N$.

\subsection{Weyl chamber as simplicial cone}

\begin{definition} A connected component of the set
$$\left\{x\in\mathbb{R}^N:\langle x,v\rangle \neq 0\ \ { for} \ { all} \ v\in R\right\}$$
is called an open Weyl chamber. Its closure is called a Weyl chamber.
\end{definition}

The following assertion can be found in e.g. Theorem 2 on p.166 in \cite{Bourbaki}.

\begin{theorem}\label{phi theorem} Let $\Gamma$ be a chamber. There exists a linear isomorphism $\Phi:\mathbb{R}^N\to\mathbb{R}^N$ such that $\Phi(\Gamma)=R_+^{N_1}\times\mathbb{R}^{N_2}$ (with $N_1+N_2=N$).
\end{theorem}




\subsection{System of dyadic cubes}

\begin{nota}\label{standardsystem}
Let
$$\mathbb{D}=\bigcup_{k\in\mathbb{Z}}\mathbb{D}_k$$
be the standard system of dyadic cubes over $\mathbb{R}^N.$ Here each dyadic cube in $\mathbb{D}_k$ is of the shape
$$\frac{m}{2^k}+[0,2^{-k})^N,\quad m\in\mathbb{Z}^N.$$

For every dyadic cube $Q\in\mathbb{D}$, we denote by $c_Q$ and $\ell(Q)$ its center and side-length, respectively.
\end{nota}


\subsection{Nearly weakly  orthonormal sequences}
Rochberg and Semmes \cite{RS} introduced the notion of \emph{nearly weakly  orthonormal (NWO)} sequences of functions. This notion has since proved highly effective in characterizing Schatten class membership of  singular integral commutators in a variety of contexts (see \cite{FLLXarxiv,HKarxiv,MR4756021,MR4552581,MR4873796,LLW,RS,Zhangwei1,Zhangwei2}). The properties of NWO sequences were originally investigated in the Euclidean setting (equipped with Lebesgue measure), but the extension to spaces of homogeneous type is straightforward once a system of dyadic cubes is available in this general framework (see \cite[Section 13]{HKarxiv} for a complete proof).

{\color{red} Definition \ref{NWO def} and Lemma \ref{NWO lemma} as stated below are my understanding of Hytonen's Corollary 8.2 and Lemma 8.7. Please check whether this understanding is the correct one!}

\begin{definition}\label{NWO def} Let $w\in A_2(\mathbb{R}^N).$ Let $c>0$ and $2<r\leq\infty.$ A sequence of functions $\{e_Q\}_{Q\in\mathbb{D}}$ is called $(r,c)$-NWO if
\begin{enumerate}
\item ${\rm supp}(e_Q)\subseteq cQ$ for every $Q\in\mathbb{D};$ 
\item $\|e_Q\|_{L_r(\mathbb{R}^N,w(x)dx)}\leq w(Q)^{\frac1r-\frac12}$ for every $Q\in\mathbb{D};$
\end{enumerate}
\end{definition}

\begin{lemma}\label{NWO lemma} Let $w\in A_2(\mathbb{R}^N).$ Let $0<p<\infty.$ Suppose that $A$ is an integral operator on $L_2(\mathbb{R}^N,w(x)dx)$ with integral kernel 
$$(x,y)\to \sum_{Q\in\mathbb{D}}\lambda_Q e_Q(x)f_Q(y),\quad x,y\in\mathbb{R}^N,$$
where $\{e_Q\}_{Q\in\mathbb{D}}$ and $\{f_Q\}_{Q\in\mathbb{D}}$ are $(r_1,c_1)$-NWO and $(r_2,c_2)$-NWO sequences, respectively. We have
$$\|A\|_{\mathcal{L}_{p,\infty}(L_2(\mathbb{R}^N,w(x)dx))}\leq C_{w,r_1,c_1,r_2,c_2,p} \|\{\lambda_Q\}_{Q\in\mathbb{D}}\|_{p,\infty},$$
$${\rm dist}_{\mathcal{L}_{p,\infty}(L_2(\Gamma,w(x)dx))}\Big(A,(\mathcal{L}_{p,\infty})_0(L_2(\mathbb{R}^N,w(x)dx))\Big)\leq $$
$$\leq C_{w,r_1,c_1,r_2,c_2,p}{\rm dist}_{l_{p,\infty}(\mathbb{D})}\Big(\{\lambda_Q\}_{Q\in\mathbb{D}},(l_{p,\infty})_0(\mathbb{D})\Big).$$
\end{lemma}
\begin{proof} The assertion follows from Corollary 8.2 and Lemma 8.7 in \cite{HKarxiv}.
\end{proof}



\section{Chamber is a doubling metric measure space}

\begin{theorem}\label{main:dunkldoubling} Let $\Gamma$ be a chamber. For all $x\in\Gamma$ and for all $r>0,$ we have
$$w(B(x,r)\cap\Gamma)\geq c_{\kappa}w(B(x,r)).$$
In particular, $(\Gamma,|\cdot|,w(x)dx)$ is a doubling metric measure space.
\end{theorem}

\begin{nota} Let $\theta:2^{N_1}\to\mathbb{R}_+.$ Set
$$w_{\theta}(x)=\prod_{A\subset\{1,\cdots,N_1\}}(\sum_{k\in A}x_k)^{\theta_A},\quad x\in\mathbb{R}_+^{N_1}.$$
\end{nota}

\begin{lemma}\label{s2 upper estimate} For every $\theta:2^{N_1}\to\mathbb{R}_+$ and for every $x\in\mathbb{R}^{N_1}_+\times\mathbb{R}^{N_2},$ we have
$$\int_{B(x,r)}w_{\theta}(y_1,\cdots,y_{N_1})dy\leq c_{N,\theta}^{(1)}r^Nw_{\theta}((x_1,\cdots,x_{N_1})+r\mathbf{1}_{N_1}).$$
\end{lemma}
\begin{proof} Obviously,
$$B(x,r)\subset \Big(\times_{k=1}^{N_1}[(x_k-r)_+,x_k+r]\Big)\times B((x_{N_1+1},\cdots,x_N),r).$$
Thus,
$$\int_{B(x,r)}w_{\theta}(y_1,\cdots,y_{N_1})dy\leq$$
$$\leq\int_{\big(\times_{k=1}^{N_1}[(x_k-r)_+,x_k+r]\big)\times B((x_{N_1+1},\cdots,x_N),r)}w_{\theta}(y_1,\cdots,y_{N_1})dy=$$
$$=\int_{\times_{k=1}^{N_1}[(x_k-r)_+,x_k+r]}w_{\theta}(u)du\times {\rm Vol}(\mathbb{B}_{N_2})r^{N_2}\leq$$
$$\leq\int_{\times_{k=1}^{N_1}[0,x_k+r]}w_{\theta}((x_1,\cdots,x_{N_1})+r\mathbf{1}_{N_1})du\times {\rm Vol}(\mathbb{B}_{N_2})r^{N_2}=$$
$$=w_{\theta}((x_1,\cdots,x_{N_1})+r\mathbf{1}_{N_1})\times\prod_{k=1}^{N_1}(x_k+r-(x_k-r)_+)du\times {\rm Vol}(\mathbb{B}_{N_2})r^{N_2}\leq $$
$$\leq w_{\theta}((x_1,\cdots,x_{N_1})+r\mathbf{1}_{N_1})\times(2r)^{N_1}\times {\rm Vol}(\mathbb{B}_{N_2})r^{N_2}.$$
\end{proof}

\begin{lemma}\label{s2 lower estimate} For every $\theta:2^{N_1}\to\mathbb{R}_+$ and for every $x\in\mathbb{R}^{N_1}_+\times\mathbb{R}^{N_2},$ we have
$$\int_{B(x,r)\cap(\mathbb{R}^{N_1}_+\times\mathbb{R}^{N_2})}w_{\theta}(y_1,\cdots,y_{N_1})dy\geq c_{N,\theta}^{(2)}r^Nw_{\theta}((x_1,\cdots,x_{N_1})+r\mathbf{1}_{N_1}).$$
\end{lemma}
\begin{proof} Obviously,
$$x+[\frac{r}{2N},\frac{r}{N}]^N\subset B(x,r)\cap (\mathbb{R}^{N_1}_+\times\mathbb{R}^{N_2}).$$
It follows that
$$\int_{B(x,r)\cap(\mathbb{R}^{N_1}_+\times\mathbb{R}^{N_2})}w_{\theta}(y_1,\cdots,y_{N_1})dy\geq $$
$$\geq\int_{x+[\frac{r}{2N},\frac{r}{N}]^N}w_{\theta}(y_1,\cdots,y_{N_1})dy=$$
$$=\int_{(x_1,\cdots,x_{N_1})+[\frac{r}{2N},\frac{r}{N}]^{N_1}}w_{\theta}(u)du\times (\frac{r}{2N})^{N_2}\geq $$
$$\geq \int_{(x_1,\cdots,x_{N_1})+[\frac{r}{2N},\frac{r}{N}]^{N_1}}w_{\theta}((x_1,\cdots,x_{N_1})+\frac{r}{2N}\mathbf{1}_{N_1})du\times (\frac{r}{2N})^{N_2}=$$
$$=w_{\theta}((x_1,\cdots,x_{N_1})+\frac{r}{2N}\mathbf{1}_{N_1})\times(\frac{r}{2N})^N\geq $$
$$\geq w_{\theta}((x_1,\cdots,x_{N_1})+r\mathbf{1}_{N_1})\times(\frac{r}{2N})^N\times (2N)^{-\sum_{A\subset\{1,\cdots,N_1\}}\theta_A}.$$
\end{proof}


\begin{lemma}\label{s2 trivial lemma} There exists $\theta:2^{N_1}\to\mathbb{R}_+$ such that
$$c_{\kappa}^{-1}w_{\theta}(x_1,\cdots,x_{N_1})\leq w(\Phi^{-1}(x))\leq c_{\kappa}w_{\theta}(x_1,\cdots,x_{N_1}),\quad x\in\mathbb{R}_+^{N_1}\times\mathbb{R}^{N_2}.$$
\end{lemma}
\begin{proof} For every $x\in\mathbb{R}^{N_1}_+\times\mathbb{R}^{N_2}$ and for every $v\in R,$ we have
$$\langle \Phi^{-1}(x),v\rangle=\langle x,(\Phi^{-1})^{\ast}(v)\rangle.$$
The latter expression has the same sign (positive or negative) on the whole cone $\mathbb{R}^{N_1}_+\times\mathbb{R}^{N_2}.$ Hence,
$$\big((\Phi^{-1})^{\ast}(v)\big)_k=
\begin{cases}
\mbox{something positive},& 1\leq k\leq N_1\\
0,& k>N_1
\end{cases}
$$
Let $A_v$ be the set of those $1\leq k\leq N_1$ with $\big((\Phi^{-1})^{\ast}(v)\big)_l$ being strictly positive. We, obviously, have
$$c_{\Phi,v}^{-1}\sum_{k\in A_v}x_k\leq |\langle \Phi^{-1}(x),v\rangle|\leq c_{\Phi,v}\sum_{k\in A_v}x_k,\quad x\in \mathbb{R}^{N_1}_+\times\mathbb{R}^{N_2}.$$
Thus,
$$\prod_{v\in R}c_{\Phi,v}^{-\kappa(v)}(\sum_{k\in A_v}x_k)^{\kappa(v)}\leq w(x)\leq \prod_{v\in R}c_{\Phi,v}^{\kappa(v)}(\sum_{k\in A_v}x_k)^{\kappa(v)},\quad x\in\mathbb{R}_+^{N_1}\times\mathbb{R}^{N_2}.$$
This completes the proof.
\end{proof}

\begin{proof}[Proof of Theorem \ref{main:dunkldoubling}] The first assertion follows by combining Lemmas \ref{s2 upper estimate}, \ref{s2 lower estimate} and \ref{s2 trivial lemma}. The second assertion is an obvious consequence of the first one.
\end{proof}


\section{First Rochberg-Semmes based estimate}

\begin{theorem}\label{first RS-type theorem} Let $N\geq 2,$ $1\leq N_1\leq N$ and $N_2=N-N_1.$ Let $\Omega=\mathbb{R}^{N_1}_+\times\mathbb{R}^{N_2}.$ Let $F$ be a face of $\Omega$ (not necessarily of maximal dimension, possibly of zero dimension). Let $\theta:2^N\to\mathbb{R}_+.$ Suppose that $K:\Omega\times\Omega\to\mathbb{C}$
satisfies
$$\bigl|\partial_x^{\alpha}\partial_y^{\beta}K\bigr|(x,y)\leq c_{K,\alpha,\beta}\frac{|x-y|^{1-|\alpha|_1-|\beta|_1}}{(|x-y|+{\rm dist}(x,F))\cdot w_{\theta}(B(x,|x-y|))}$$
for every $x,y\in\Omega$ and for every $\alpha,\beta\in\mathbb{Z}^N_+.$ Let $V$ be an integral operator on $L_2(\Omega,w_{\theta}(x)dx)$ with integral kernel $K.$ For every $f\in C^{\infty}_c(\mathbb{R}^N)$ vanishing on $F,$ we have
$$\|M_fV\|_{\mathcal{L}_{N,\infty}(L_2(\Omega,w_{\theta}(x)dx))}\leq c_{K,\theta}\||\nabla f|\|_{L_N(\Omega)}.$$
\end{theorem}

\begin{theorem}\label{first RS-type convergence} Suppose we are the setting of Theorem \ref{first RS-type theorem}. If $0\leq\phi_m\leq 1$ for $m\geq1$ are in $C^{\infty}_c(\mathbb{R}^N)$ and ${\rm supp}(1-\phi_m)$ increases to $\mathbb{R}^N,$ then
$${\rm dist}_{\mathcal{L}_{N,\infty}(L_2(\Omega,w_{\theta}(x)dx))}\Big(M_{\phi_m}\cdot M_fV\cdot M_{\phi_m}-M_fV,(\mathcal{L}_{N,\infty})_0(L_2(\Omega,w_{\theta}(x)dx))\Big)\to0$$
as $m\to\infty.$
\end{theorem}

\subsection{Proof of Theorem \ref{first RS-type theorem}}


\begin{lemma} There exists a smooth function $\psi$ supported on $[-1,1]^N\backslash(-\frac14,\frac14)^N$ such that
$$1=\sum_{Q\in\mathbb{D}}\chi_Q(x)\chi_{3Q}(y)\psi(\frac{x-y}{\ell(Q)}).$$
\end{lemma}
\begin{proof} Choose $\rho\in C_c^{\infty}(\mathbb{R}^N)$ such that $\rho=1$ on $(-\frac12,\frac12)^N$ and $\rho=0$ outside $(-1,1)^N.$ Set $\psi(z)=\rho(z)-\rho(2z).$ Clearly,
$${\rm supp}(\psi)\subset
\left\{z\in\mathbb{R}^N:\quad 
\frac14\leq |z|_{\infty}\leq 1\right\}.$$
For every $z\neq0,$ we have
\begin{equation}\label{eq:rankone-psi-partition}
\sum_{j\in\mathbb{Z}}\psi(2^jz)=1.
\end{equation}

For every $x,$ there exists a unique $Q\in\mathbb{D}_j$ containing $x.$ For this $Q,$ we have $\ell(Q)=2^{-j}.$ We have
$$1=\sum_{j\in\mathbb{Z}}\psi(2^j(x-y))=\sum_{j\in\mathbb{Z}}(\sum_{Q\in\mathbb{D}_j}\chi_Q(x))\psi(2^j(x-y))=$$
$$=\sum_{j\in\mathbb{Z}}\sum_{Q\in\mathbb{D}_j}\chi_Q(x)\psi(\frac{x-y}{\ell(Q)})=\sum_{Q\in\mathbb{D}}\chi_Q(x)\psi(\frac{x-y}{\ell(Q)}).$$

If $x\in Q$ and $\psi(\frac{x-y}{\ell_Q})\neq 0,$ then
$$x\in Q,\quad |x-y|_{\infty}\leq\ell(Q).$$
We obviously have $|x-c_Q|_{\infty}\leq\frac12\ell(Q)$ and $|x-y|_{\infty}\leq\ell(Q).$ Thus, $|y-c_Q|_{\infty}\leq\frac32\ell(Q).$ In other words, $y\in 3Q.$ Thus,
$$\chi_Q(x)\psi(\frac{x-y}{\ell(Q)})=\chi_Q(x)\chi_{3Q}(y)\psi(\frac{x-y}{\ell(Q)}).$$
This completes the proof.
\end{proof}

\begin{lemma}\label{lem:Weyl-off-diagonal-kernel-NWO} In the setting of Theorem \ref{first RS-type theorem}, we have
$$f(x)K(x,y)=\sum_{k,l\in\mathbb{Z}^N}(1+|k|^2+|l|^2)^{-N-1}\sum_{\substack{Q\in\mathbb{D}\\ Q\subset\Omega}}e_{k,l,Q}(x)f_{k,l,Q}(y),$$
where
\begin{enumerate}[{\rm (i)}]
\item $e_{k,l,Q}$ is supported in $Q$ and $f_{k,l,Q}$ is supported in $3Q$ for every $Q\in\mathbb{D}$ with $Q\subset\Omega;$
\item for every $k,l\in\mathbb{Z}^N$ and for every $Q\in\mathbb{D}$ with $Q\subset\Omega,$ we have
$$\|e_{k,l,Q}\|_{L_{2N}(\mathbb{R}^N,w_{\theta}(x)dx)}=\|f\|_{L_{2N}(Q,w_{\theta}(x)dx)};$$
\item for every $k,l\in\mathbb{Z}^N$ and for every $Q\in\mathbb{D}$ with $Q\subset\Omega,$ we have
$$\|f_{k,l,Q}\|_{L_{\infty}(\mathbb{R}^N)}\leq c_{K,\kappa}\frac{\ell(Q)}{{\rm dist}(c_Q,F)};$$
\end{enumerate}
\end{lemma}
\begin{proof} For every $x,y\in \Omega,$ we have
$$K(x,y)=\sum_{Q\in\mathbb{D}}\chi_Q(x)\chi_{3Q}(y)K(x,y)\psi(\frac{x-y}{\ell(Q)})=$$
$$=\sum_{\substack{Q\in\mathbb{D}\\ Q\subset\Omega}}\chi_Q(x)\chi_{3Q}(y)K(x,y)\psi(\frac{x-y}{\ell(Q)})=$$

For every $Q\in\mathbb{D}$ with $3Q\subset\Omega,$ define the map
$$K_Q:[-\frac12,\frac12]^N\times[-\frac32,\frac32]^N\to\mathbb{C}$$
by setting
$$K_Q(a,b)=K(c_Q+a\ell(Q),c_Q+b\ell(Q))\psi(a-b),\quad a\in [-\frac12,\frac12]^N,\quad b\in[-\frac32,\frac32]^N.$$
If $Q\in\mathbb{D}$ is such that $3Q\not\subset\Omega,$ then the definition of $K_Q$ is the same, $a$ also runs through $[-\frac12,\frac12]^N,$ but the parallelepiped $b$ runs through is slightly smaller. {\color{red} to co-authors: if you wish, specify the domain of $b$ also.} {\it In what follows, we lighten the notations by pretending that $b$ is always defined on $[-\frac32,\frac32]^N.$}

It follows from the assumption on $K$ that
$$\|K_Q\|_{C^{2N+2}([-\frac32,\frac32]^N\times[-\frac32,\frac32]^N)}\leq\frac{c_{K,\theta}'}{w_{\theta}(Q)}\cdot \frac{\ell(Q)}{{\rm dist}(c_Q,F)}.$$
Hence, there exists an extension $K_Q:\mathbb{T}^N\times\mathbb{T}^N\to\mathbb{C}$ such that
$$\|K_Q\|_{C^{2N+2}(\mathbb{T}^N\times\mathbb{T}^N)}\leq \frac{c_{K,\theta}''}{w_{\theta}(Q)}\cdot \frac{\ell(Q)}{{\rm dist}(c_Q,F)}.$$
In particular,
$$\|K_Q\|_{W^{2N+2.2}(\mathbb{T}^N\times\mathbb{T}^N)}\leq \frac{c_{K,\theta}}{w_{\theta}(Q)}\cdot \frac{\ell(Q)}{{\rm dist}(c_Q,F)}.$$
We may, therefore, write
$$K_Q(a,b)=\sum_{k,l\in\mathbb{Z}^N}{\rm coef}_{k,l,Q}e^{i\langle k,a\rangle}e^{i\langle l,b\rangle},\quad a,b\in\mathbb{T}^N,$$
where
$$\sum_{k,l\in\mathbb{Z}^N}(1+|k|^2+|l|^2)^{2N+2}|{\rm coef}_{k,l,Q}|^2\leq \Big(\frac{c_{K,\theta}}{w_{\theta}(Q)}\cdot \frac{\ell(Q)}{{\rm dist}(c_Q,F)}\Big)^2.$$
In particular,
$$(1+|k|^2+|l|^2)^{N+1}\cdot|{\rm coef}_{k,l,Q}|\leq \frac{c_{K,\theta}}{w_{\theta}(Q)}\cdot \frac{\ell(Q)}{{\rm dist}(c_Q,F)},\quad k,l\in\mathbb{Z}^N.$$

We now set
$$e_{k,l,Q}(x)=f(x)\chi_Q(x)e^{i\langle k,\frac{x-c_Q}{\ell(Q)}\rangle},$$
$$f_{k,l,Q}(y)=(1+|k|^2+|l|^2)^{N+1}{\rm coef}_{k,l,Q}\cdot \chi_{3Q}(y)e^{i\langle l,\frac{y-c_Q}{\ell(Q)}\rangle}.$$
This allows us to write
$$f(x)K(x,y)=\sum_{k,l\in\mathbb{Z}^N}(1+|k|^2+|l|^2)^{-N-1}\sum_{\substack{Q\in\mathbb{D}\\ Q\subset\Omega}}e_{k,l,Q}(x)f_{k,l,Q}(y).$$

Obviously,
$${\rm supp}(e_{k,l,Q})\subset Q,\quad {\rm supp}(f_{k,l,Q})\subset 3Q,$$
$$\|e_{k,l,Q}\|_{L_{2N}(\mathbb{R}^N,w_{\theta}(x)dx)}=\|f\|_{L_{2N}(Q,w(x)dx)},$$
$$\|f_{k,l,Q}\|_{L_{\infty}(\mathbb{R}^N)}=(1+|k|^2+|l|^2)^{N+1}|{\rm coef}_{k,l,Q}|\leq \frac{c_{K,\theta}}{w_{\theta}(Q)}\cdot \frac{\ell(Q)}{{\rm dist}(c_Q,F)}.$$
\end{proof}

\begin{lemma}\label{supremum vs average} For every $Q\in\mathbb{D}$ with $Q\subset\Omega,$ we have
$$\sup_{x\in Q}w_{\theta}(x)\leq c_{\theta}\frac{w_{\theta}(Q)}{m(Q)}.$$
\end{lemma}
\begin{proof} We have
$$w(Q)\geq c_{\theta}\prod_{A\subset\{1,\cdots,N\}}(\ell_Q+\sum_{k\in A}(c_Q)_k)^{\theta_A}.$$
On the other hand, if $x\in Q,$ then $|x-c_Q|_{\infty}\leq \frac12\ell(Q)$ and, therefore,
$$\sum_{k\in A}x_k\leq \frac{|A|}{2}\ell_Q+\sum_{k\in A}(c_Q)_k\leq N(\ell_Q+\sum_{k\in A}(c_Q)_k).$$
Thus,
$$w(x)=\prod_{A\subset\{1,\cdots,N\}}(\sum_{k\in A}x_k)^{\theta_A}\leq N^{\sum_{A\subset\{1,\cdots,N\}}\theta_A}\prod_{A\subset\{1,\cdots,N\}}(\ell_Q+\sum_{k\in A}(c_Q)_k)^{\theta_A}.$$
This completes the proof.
\end{proof}

\begin{lemma}\label{first RS intermediate estimate} In the setting of Theorem \ref{first RS-type theorem}, we have
$$\|M_fV\|_{\mathcal{L}_{N,\infty}(L_2(\Omega,w_{\theta}(x)dx))}\leq c_{K,\theta}\Big\|\Big\{m(Q)^{-\frac1{2N}}\|f\|_{L_{2N}(Q)}\cdot\frac{\ell(Q)}{{\rm dist}(c_Q,F)}\Big\}_{\substack{Q\in\mathbb{D}\\ Q\subset\Omega}}
\Big\|_{N,\infty}.$$
\end{lemma}
\begin{proof} Let $A_{k,l}$ be the integral operator with integral kernel given by the formula
$$(x,y)\to\sum_{\substack{Q\in\mathbb{D}\\ Q\subset\Omega}}e_{k,l,Q}(x)f_{k,l,Q}(y),\quad x,y\in\Omega,\quad k,l\in\mathbb{Z}^N.$$
By Lemma \ref{lem:Weyl-off-diagonal-kernel-NWO}, we have
$$M_fV=\sum_{k,l\in\mathbb{Z}^N}(1+|k|^2+|l|^2)^{-N-1}A_{k,l}.$$
By the triangle inequality,
$$\|M_fV\|_{\mathcal{L}_{N,\infty}(L_2(\Omega,w_{\theta}(x)dx))}\leq$$
$$\leq \frac{N}{N-1}\sum_{k,l\in\mathbb{Z}^N}(1+|k|^2+|l|^2)^{-N-1} \|A_{k,l}\|_{\mathcal{L}_{N,\infty}(L_2(\Omega,w_{\theta}(x)dx))}.$$

By Lemma \ref{lem:Weyl-off-diagonal-kernel-NWO}, the sequences
$$\Big\{w_{\theta}(Q)^{\frac1{2N}-\frac12}\frac{e_{k,l,Q}}{\|f\|_{L_{2N}(Q,w_{\theta}(x)dx)}}\Big\}_{\substack{Q\in\mathbb{D}\\ Q\subset\Omega}},\quad \Big\{c_{K,\theta}^{-1}w_{\theta}(Q)^{\frac12}\cdot\frac{{\rm dist}(c_Q,F)}{\ell_Q}f_{k,l,Q}\Big\}_{\substack{Q\in\mathbb{D}\\ Q\subset\Omega}}$$
are $(2N,1)$-NWO and, respectively, $(\infty,3)$-NWO in the sense of Definition \ref{NWO def}. Using Lemma \ref{NWO lemma}, we conclude that
$$\|A_{k,l}\|_{\mathcal{L}_{N,\infty}(L_2(\Omega,w_{\theta}(x)dx))}\leq$$
$$\leq c_{\theta,2N,1,\infty,3,N}c_{K,\theta}\Big\|\Big\{w_{\theta}(Q)^{-\frac1{2N}}\|f\|_{L_{2N}(Q,w_{\theta}(x)dx)}\cdot\frac{\ell_Q}{{\rm dist}(c_Q,F)}\Big\}_{\substack{Q\in\mathbb{D}\\ Q\subset\Omega}}\Big\|_{N,\infty}.$$
Therefore,
$$\|M_fV\|_{\mathcal{L}_{N,\infty}(L_2(\mathbb{R}^N_+,w_{\theta}(x)dx))}\leq$$
$$\leq \frac{N}{N-1}c_{\theta,2N,1,\infty,3,N}c_{K,\theta}\sum_{k,l\in\mathbb{Z}^N}(1+|k|^2+|l|^2)^{-N-1}\times$$
$$\times \Big\|\Big\{w_{\theta}(Q)^{-\frac1{2N}}\|f\|_{L_{2N}(Q,w_{\theta}(x)dx)}\cdot\frac{\ell_Q}{{\rm dist}(c_Q,F)}\Big\}_{\substack{Q\in\mathbb{D}\\ Q\subset\Omega}}\Big\|_{N,\infty}.$$
By Lemma \ref{supremum vs average},
$$w_{\theta}(Q)^{-\frac1{2N}}\|f\|_{L_{2N}(Q,w_{\theta}(x)dx)}\leq c_{\theta}^{\frac1{2N}}m(Q)^{-\frac1{2N}}\|f\|_{L_{2N}(Q)},\quad Q\in\mathbb{D},\quad Q\subset\Omega.$$
Combining the last two estimates, we complete the proof.
\end{proof}

\begin{lemma}\label{first RS prefinal lemma} Let $\Omega$ and $F$ be as in Theorem \ref{first RS-type theorem}. For every $f\in C^{\infty}_c(\mathbb{R}^N)$ vanishing on $F,$ we have
$$\Big\|\Big\{m(Q)^{-1}\|f\|_{L_1(Q)}\cdot\frac{\ell(Q)}{{\rm dist}(c_Q,F)}\Big\}_{\substack{Q\in\mathbb{D}\\ Q\subset\Omega}}
\Big\|_{N,\infty}\leq c_N\||\nabla f|\|_{L_N(\Omega)}.$$
{\color{red} I proved the lemma for the case $F\neq\{0\}.$}
\end{lemma}
\begin{proof} Choose $S\subset\{1,\cdots,N_1\}$ such that
$$F=\{x\in\mathbb{R}^N:\ x_k\geq0,\ k\in S,\ x_l=0,\ 1\leq l\leq N_1,\ l\notin S\}.$$
Define dilation $D_t:\mathbb{R}^N\to\mathbb{R}^N$ by setting
$$(D_tx)_j=
\begin{cases}
tx_j,& 1\leq j\leq N_1,\ j\notin S\\
x_j,&\mbox{otherwise}
\end{cases}
$$
We obviously have
$$f(x)=f(D_1x)-f(D_0x)=\int_0^1\frac{d}{dt}f(D_tx)dt=\sum_{\substack{1\leq j\leq N_1\\ j\notin S}}x_j\cdot \int_0^1(\partial_j f)(D_tx)dt.$$
Since $x_j\leq{\rm dist}(x,F),$ it follows that
$$|f(x)|\leq {\rm dist}(x,F)\cdot \sum_{j=1}^N\int_0^1|(\partial_jf)|(D_tx)dt.$$
In other words,
$$|f|\leq {\rm dist}(x,F)\cdot \sum_{j=1}^N\int_0^1(|\partial_jf|\circ D_t)dt\leq {\rm dist}(x,F)\cdot N\int_0^1(|\nabla f|\circ D_t)dt.$$

If $x\in Q,$ then
$${\rm dist}(x,F)\leq{\rm dist}(c_Q,F)+|x-c_Q|\leq {\rm dist}(c_Q,F)+\frac{N}{2}\ell(Q)\leq (N+1){\rm dist}(c_Q,F).$$
Therefore,
$$\frac1{{\rm dist}(c_Q,F)}\|f\|_{L_1(Q)}\leq N(N+1)\Big\|\int_0^1(|\nabla f|\circ D_t)dt\Big\|_{L_1(Q)}.$$
Thus,
$$\Big\|\Big\{m(Q)^{-1}\|f\|_{L_1(Q)}\cdot\frac{\ell(Q)}{{\rm dist}(c_Q,F)}\Big\}_{\substack{Q\in\mathbb{D}\\ Q\subset\Omega}}
\Big\|_{N,\infty}\leq $$
$$\leq N(N+1)\Big\|\Big\{m(Q)^{\frac1N-1}\Big\|\int_0^1(|\nabla f|\circ D_t)dt\Big\|_{L_1(Q)}\Big\}_{\substack{Q\in\mathbb{D}\\ Q\subset\Omega}}
\Big\|_{N,\infty}.$$
Using Lemma \ref{lem:Cwikel-dyadic-distribution}, we write
$$\Big\|\Big\{m(Q)^{-1}\|f\|_{L_1(Q)}\cdot\frac{\ell(Q)}{{\rm dist}(c_Q,F)}\Big\}_{\substack{Q\in\mathbb{D}\\ Q\subset\Omega}}
\Big\|_{N,\infty}\leq $$
$$\leq C_N\Big\|\int_0^1(|\nabla f|\circ D_t)dt\Big\|_{L_N(\Omega)}\leq$$
$$\leq C_N\int_0^1\||\nabla f|\circ D_t\|_{L_N(\Omega)}dt=C_N\||\nabla f|\|_{L_N(\Omega)}\cdot\int_0^1t^{\frac{|S|-N_1}{N}}dt.$$
Since $|S|-N_1>-N,$ the assertion follows.
\end{proof}


\begin{lemma}\label{first RS final lemma} Let $\Omega$ and $F$ be as in Theorem \ref{first RS-type theorem}. For every $f\in C^{\infty}_c(\mathbb{R}^N)$ vanishing on $F,$ we have
$$\Big\|\Big\{m(Q)^{-\frac1{2N}}\|f\|_{L_{2N}(Q)}\cdot\frac{\ell(Q)}{{\rm dist}(c_Q,F)}\Big\}_{\substack{Q\in\mathbb{D}\\ Q\subset\Omega}}
\Big\|_{N,\infty}\leq c_N\||\nabla f|\|_{L_N(\mathbb{R}^N)}.$$
\end{lemma}
\begin{proof} We obviously have
$$\|f\|_{L_{2N}(Q)}\leq\Big\|f-\frac1{m(Q)}\int_Qf\Big\|_{L_{2N}(Q)}+m(Q)^{\frac1{2N}-1}\|f\|_{L_1(Q)}.$$
By Sobolev inequality,
$$\|f\|_{L_{2N}(Q)}\leq c_N\|\nabla f\|_{L_{\frac{2N}{3}}(Q)}+m(Q)^{\frac1{2N}-1}\|f\|_{L_1(Q)}.$$
By the triangle inequality,
$$\Big\|\Big\{m(Q)^{-\frac1{2N}}\|f\|_{L_{2N}(Q)}\cdot\frac{\ell(Q)}{{\rm dist}(c_Q,F)}\Big\}_{\substack{Q\in\mathbb{D}\\ Q\subset\Omega}}
\Big\|_{N,\infty}\leq$$
$$\leq\frac{2N}{N-1}c_N\Big\|\Big\{m(Q)^{-\frac1{2N}}\|\nabla f\|_{L_{\frac{2N}{3}}(Q)}\Big\}_{\substack{Q\in\mathbb{D}\\ Q\subset\Omega}}\Big\|_{N,\infty}+$$
$$+\frac{N}{N-1}\Big\|\Big\{m(Q)^{-1}\|f\|_{L_1(Q)}\cdot\frac{\ell(Q)}{{\rm dist}(c_Q,F)}\Big\}_{\substack{Q\in\mathbb{D}\\ Q\subset\Omega}}
\Big\|_{N,\infty}.$$
The assertion follows now from Lemma \ref{lem:Cwikel-dyadic-distribution} (for the first summand) and Lemma \ref{first RS prefinal lemma} (for the second summand).
\end{proof}


\begin{proof}[Proof of Theorem \ref{first RS-type theorem}] The assertion follows by combining Lemmas \ref{first RS intermediate estimate} and \ref{first RS final lemma}.
\end{proof}

\subsection{Proof of Theorem \ref{first RS-type convergence}}

\begin{lemma}\label{first RS convergence lemma} Let $f,\phi\in C^{\infty}_c(\mathbb{R}^N).$ If $0\leq\phi\leq1,$ then
$${\rm dist}_{\mathcal{L}_{N,\infty}(L_2(\Gamma,w(x)dx))}\Big(M_fV-M_{f\phi}V,(\mathcal{L}_{N,\infty})_0(L_2(\Gamma,w(x)dx))\Big)\leq c_{K,\kappa,\phi,r}\times$$
$$\times{\rm dist}_{l_{N,\infty}(\mathbb{D}_{\Gamma})}\Big(\Big\{\frac{d_Im(I)^{-\frac1r}}{|c_I-gc_I|}\|f\|_{L_r(I)}\Big\}_{\substack{I\in\mathbb{D}_{\Gamma}\\ I\cap {\rm supp}(f)\neq\emptyset\\ I\cap {\rm supp}(1-\phi)\neq\emptyset}},(l_{N,\infty})_0(\mathbb{D}_{\Gamma})\Big),$$
$${\rm dist}_{\mathcal{L}_{N,\infty}(L_2(\Gamma,w(x)dx))}\Big(M_fV-M_fVM_{\phi},(\mathcal{L}_{N,\infty})_0(L_2(\Gamma,w(x)dx))\Big)\leq c_{K,\kappa,\phi,r}\times$$
$$\times{\rm dist}_{l_{N,\infty}(\mathbb{D}_{\Gamma})}\Big(\Big\{\frac{d_Im(I)^{-\frac1r}}{|c_I-gc_I|}\|f\|_{L_r(I)}\Big\}_{\substack{I\in\mathbb{D}_{\Gamma}\\ I\cap {\rm supp}(f)\neq\emptyset\\ 3I\cap {\rm supp}(1-\phi)\neq\emptyset}},(l_{N,\infty})_0(\mathbb{D}_{\Gamma})\Big).$$
\end{lemma}
\begin{proof} We only prove the second inequality. The proof of the first one is similar, but simpler.

Let $A_{k,l}$ be as in the proof of Lemma \ref{first RS intermediate estimate}. We have
$$M_fV-M_fVM_{\phi}=\sum_{k,l\in\mathbb{Z}^N}(1+|k|^2+|l|^2)^{-N-1}A_{k,l}M_{1-\phi}.$$
By the triangle inequality,
$${\rm dist}_{\mathcal{L}_{N,\infty}(L_2(\Gamma,w(x)dx))}\Big(M_fV-M_fVM_{\phi},(\mathcal{L}_{N,\infty})_0(L_2(\Gamma,w(x)dx))\Big)\leq$$
$$\leq \frac{N}{N-1}\sum_{k,l\in\mathbb{Z}^N}\Big((1+|k|^2+|l|^2)^{-N-1}\times$$
$$\times{\rm dist}_{\mathcal{L}_{N,\infty}(L_2(\Gamma,w(x)dx))}\Big(A_{k,l}M_{1-\phi},(\mathcal{L}_{N,\infty})_0(L_2(\Gamma,w(x)dx))\Big)\Big).$$


By Lemma \ref{lem:Weyl-off-diagonal-kernel-NWO}, the sequences
$$\Big\{w(I)^{\frac1r-\frac12}\frac{e_{k,l,I}}{\|f\|_{L_r(I,w(x)dx)}}\Big\}_{I\in\mathbb{D}_{\Gamma}},\quad \Big\{c_{K,\kappa,\Phi}^{-1}w(I)^{\frac12}\cdot\frac{|c_I-gc_I|}{d_I}f_{k,l,I}\cdot(1-\phi)\Big\}_{I\in\mathbb{D}_{\Gamma}}$$
are $(r,1)$-NWO and, respectively, $(\infty,3)$-NWO in the sense of Definition \ref{NWO Gamma def}. Obviously, $e_{k,l,I}\not\equiv0$ only if $I\cap {\rm supp}(f)\neq\emptyset$ and $f_{k,l,I}\cdot(1-\phi)\not\equiv 0$ only if $I\cap {\rm supp}(f)\neq\emptyset$ and $3I\cap {\rm supp}(1-\phi)\neq\emptyset.$ Using Lemma \ref{NWO Gamma lemma}, we conclude that
$${\rm dist}_{\mathcal{L}_{N,\infty}(L_2(\Gamma,w(x)dx))}\Big(A_{k,l}(1-\phi),(\mathcal{L}_{N,\infty})_0(L_2(\Gamma,w(x)dx))\Big)\leq$$
$$\leq c_{\Phi,\kappa,r,1,\infty,3,N}c_{K,\kappa,\Phi}\times$$
$$\times{\rm dist}_{l_{N,\infty}(\mathbb{D}_{\Gamma})}\Big(\Big\{\frac{d_Iw(I)^{-\frac1r}}{|c_I-gc_I|}\|f\|_{L_r(I,w(x)dx)}\Big\}_{\substack{I\in\mathbb{D}_{\Gamma}\\ I\cap {\rm supp}(f)\neq\emptyset\\ 3I\cap {\rm supp}(1-\phi)\neq\emptyset}},(l_{N,\infty})_0(\mathbb{D}_{\Gamma})\Big).$$
By Lemma \ref{supremum vs average},
$$w(I)^{-\frac1r}\|f\|_{L_r(I,w(x)dx)}\leq c_{\kappa,\Phi}^{\frac1r}m(I)^{-\frac1r}\|f\|_{L_r(I)},\quad I\in\mathbb{D}_{\Gamma}.$$
Thus,
$${\rm dist}_{\mathcal{L}_{N,\infty}(L_2(\Gamma,w(x)dx))}\Big(A_{k,l}(1-\phi),(\mathcal{L}_{N,\infty})_0(L_2(\Gamma,w(x)dx))\Big)\leq$$
$$\leq c_{\Phi,\kappa,r,1,\infty,3,N}c_{K,\kappa,\Phi}c_{\kappa,\Phi}^{\frac1r}\times$$
$$\times{\rm dist}_{l_{N,\infty}(\mathbb{D}_{\Gamma})}\Big(\Big\{\frac{d_Im(I)^{-\frac1r}}{|c_I-gc_I|}\|f\|_{L_r(I)}\Big\}_{\substack{I\in\mathbb{D}_{\Gamma}\\ I\cap {\rm supp}(f)\neq\emptyset\\ 3I\cap {\rm supp}(1-\phi)\neq\emptyset}},(l_{N,\infty})_0(\mathbb{D}_{\Gamma})\Big).$$
Combining this estimate with the one in the first paragraph {\color{red} assign label and refer to it} yields the second inequality.
\end{proof}

\begin{lemma}\label{first RS convergence for sets} For a bounded Borel $A\subset\mathbb{R}^N,$ we have
$$\lim_{n\to\infty}\Big\|\Big\{m(Q)^{\frac1N}\Big\}_{\substack{Q\in\mathbb{D}\\ Q\subset A\\ \ell(Q)\leq 2^{-n}}}\Big\|_{l_{N,\infty}(\mathbb{D})}\leq c_Nm(A)^{\frac1N}.$$
\end{lemma}
\begin{proof} For a given $\epsilon>0,$ choose $j$ (depending on $\epsilon$) and finite set $B\subset\mathbb{D}_j$ such that $A\subset \cup_{Q'\in B}Q'$ and such that $m((\cup_{Q'\in B}Q')\backslash A)<\epsilon.$ For $n\geq m,$ we have (here, we use the notation $Q_0=[0,1)^N$)
$$\Big\|\Big\{m(Q)^{\frac1N}\Big\}_{\substack{Q\in\mathbb{D}\\ Q\subset A\\ \ell(Q)\leq 2^{-n}}}\Big\|_{l_{N,\infty}(\mathbb{D})}\leq\Big\|\Big\{m(Q)^{\frac1N}\Big\}_{\substack{Q\in\mathbb{D}\\ Q\subset Q'\mbox{ for some }Q'\in B}}\Big\|_{l_{N,\infty}(\mathbb{D})}=$$
$$=(\frac{|B|}{2^j})^{\frac1N}\Big\|\Big\{m(Q)^{\frac1N}\Big\}_{\substack{Q\in\mathbb{D}\\ Q\subset Q_0}}\Big\|_{l_{N,\infty}(\mathbb{D})}=c_N\cdot m^{\frac1N}(\cup_{Q'\in B}Q')\leq c_N\cdot (m(A)+\epsilon)^{\frac1N}.$$
Passing $n\to\infty$ and, after that, passing $\epsilon\to0,$ we complete the proof.
\end{proof}

\begin{lemma}\label{first RS extra convergence for sets} We have
$$\lim_{n\to\infty}\Big\|\Big\{m(I)^{\frac1N}\Big\}_{\substack{I\in\mathbb{D}_{\Gamma}\\ I\cap {\rm supp}(f)\neq\emptyset\\ 3I\cap {\rm supp}(1-\phi)\neq\emptyset\\ d_I\leq 2^{-n}}}\Big\|_{l_{N,\infty}(\mathbb{D}_{\Gamma})}\leq c_Nm^{\frac1N}({\rm supp}(f)\cap{\rm supp}(1-\phi)).$$
\end{lemma}
\begin{proof} Denote for brevity $A=\Phi^{-1}({\rm supp}(f))$ and $B=\Phi^{-1}({\rm supp}(1-\phi)).$ For $\epsilon>0,$ denote by $A_{\epsilon}$ (respectively, $B_{\epsilon}$) the $\epsilon$-neighborhood of $A$ (respectively, of $B$) with respect to the norm $|\cdot|_{\infty}.$

If $2^{-n}<\epsilon,$ then the condition $I\cap{\rm supp}(\phi)\neq\emptyset$ yields, for $Q=\Phi(I)\in\mathbb{D},$ the condition $Q\subset A_{\epsilon}.$ If $2^{-n}<\frac{\epsilon}{3},$ then the condition $I\cap{\rm supp}(1-\phi)\neq\emptyset$ yields, for $Q=\Phi(I)\in\mathbb{D},$ the condition $Q\subset B_{\epsilon}.$ Hence, for $2^{-n}<\frac{\epsilon}{3},$ we have
$$\Big\|\Big\{m(I)^{\frac1N}\Big\}_{\substack{I\in\mathbb{D}_{\Gamma}\\ I\cap {\rm supp}(f)\neq\emptyset\\ 3I\cap {\rm supp}(1-\phi)\neq\emptyset\\ d_I\leq 2^{-n}}}\Big\|_{l_{N,\infty}(\mathbb{D}_{\Gamma})}\leq{\rm det}(\Phi)^{-\frac1N} \Big\|\Big\{m(Q)^{\frac1N}\Big\}_{\substack{Q\in\mathbb{D}\\ Q\subset A_{\epsilon}\cap B_{\epsilon}\\ \ell(Q)\leq 2^{-n}}}\Big\|_{l_{N,\infty}(\mathbb{D})}.$$
Using Lemma \ref{first RS convergence for sets} and passing $n\to\infty,$ we obtain
$$\lim_{n\to\infty}\Big\|\Big\{m(I)^{\frac1N}\Big\}_{\substack{I\in\mathbb{D}_{\Gamma}\\ I\cap {\rm supp}(f)\neq\emptyset\\ 3I\cap {\rm supp}(1-\phi)\neq\emptyset\\ d_I\leq 2^{-n}}}\Big\|_{l_{N,\infty}(\mathbb{D}_{\Gamma})}\leq c_N{\rm det}(\Phi)^{-\frac1N}m^{\frac1N}(A_{\epsilon}\cap B_{\epsilon}).$$
Since $A$ and $B$ are both closed, it follows that
$$\lim_{\epsilon\to0}m(A_{\epsilon}\cap B_{\epsilon})=m(A\cap B).$$
Passing $\epsilon\to0$ in the second last display, we complete the proof.
\end{proof}

\begin{lemma}\label{first RS with r convergence lemma} Let $r>2.$ For every $f\in C^{\infty}_c(\mathbb{R}^N)$ and for every $n\in\mathbb{N},$ we have
$${\rm dist}_{l_{N,\infty}(\mathbb{D}_{\Gamma})}\Big(\Big\{m(I)^{\frac1N-\frac1r}\|\nabla f\|_{L_r(I)}\Big\}_{\substack{I\in\mathbb{D}_{\Gamma}\\ I\cap {\rm supp}(f)\neq\emptyset\\ 3I\cap {\rm supp}(1-\phi)\neq\emptyset}},(l_{N,\infty})_0(\mathbb{D}_{\Gamma})\Big)\leq$$
$$\leq c_N\|\nabla f\|_{L_{\infty}(\mathbb{R}^N)}\times m^{\frac1N}({\rm supp}(f)\cap{\rm supp}(1-\phi)).$$
\end{lemma}
\begin{proof} For a given $n\in\mathbb{Z},$ we have
$$\Big\{m(I)^{\frac1N-\frac1r}\|\nabla f\|_{L_r(I)}\Big\}_{\substack{I\in\mathbb{D}_{\Gamma}\\ I\cap {\rm supp}(f)\neq\emptyset\\ 3I\cap {\rm supp}(1-\phi)\neq\emptyset\\ d_I>2^{-n}}}\leq\|\nabla f\|_{L_r(\mathbb{R}^N)}\Big\{m(I)^{\frac1N-\frac1r}\Big\}_{\substack{I\in\mathbb{D}_{\Gamma}\\ I\cap {\rm supp}(f)\neq\emptyset\\ 3I\cap {\rm supp}(1-\phi)\neq\emptyset\\ d_I>2^{-n}}}.$$
The right hand side obviously belongs to $(l_{N,\infty})_0(\mathbb{D}_{\Gamma})$ and, hence, so does the left hand side. Thus,
$${\rm dist}_{l_{N,\infty}(\mathbb{D}_{\Gamma})}\Big(\Big\{m(I)^{\frac1N-\frac1r}\|\nabla f\|_{L_r(I)}\Big\}_{\substack{I\in\mathbb{D}_{\Gamma}\\ I\cap {\rm supp}(f)\neq\emptyset\\ 3I\cap {\rm supp}(1-\phi)\neq\emptyset}},(l_{N,\infty})_0(\mathbb{D}_{\Gamma})\Big)\leq$$
$$\leq\Big\|\Big\{m(I)^{\frac1N-\frac1r}\|\nabla f\|_{L_r(I)}\Big\}_{\substack{I\in\mathbb{D}_{\Gamma}\\ I\cap {\rm supp}(f)\neq\emptyset\\ 3I\cap {\rm supp}(1-\phi)\neq\emptyset\\ d_I\leq 2^{-n}}}\Big\|_{l_{N,\infty}(\mathbb{D}_{\Gamma})}\leq$$
$$\leq\|\nabla f\|_{L_{\infty}(\mathbb{R}^N)}\times\Big\|\Big\{m(I)^{\frac1N}\Big\}_{\substack{I\in\mathbb{D}_{\Gamma}\\ I\cap {\rm supp}(f)\neq\emptyset\\ 3I\cap {\rm supp}(1-\phi)\neq\emptyset\\ d_I\leq 2^{-n}}}\Big\|_{l_{N,\infty}(\mathbb{D}_{\Gamma})}.$$
Passing $n\to\infty$ and using Lemma \ref{first RS extra convergence for sets}, we complete the proof.
\end{proof}

\begin{lemma}\label{first RS without r convergence lemma} For every $f\in C^{\infty}_c(\mathbb{R}^N)$ and for every $n\in\mathbb{N},$ we have
$${\rm dist}_{l_{N,\infty}(\mathbb{D}_{\Gamma})}\Big(\Big\{\frac{d_Im(I)^{-1}}{|c_I-gc_I|}\|f-f\circ g\|_{L_1(I)}\Big\}_{\substack{I\in\mathbb{D}_{\Gamma}\\ I\cap {\rm supp}(f)\neq\emptyset\\ 3I\cap {\rm supp}(1-\phi)\neq\emptyset}},(l_{N,\infty})_0(\mathbb{D}_{\Gamma})\Big)\leq$$
$$\leq c_{\Phi}\|\nabla f\|_{L_{\infty}(\mathbb{R}^N)}\times m^{\frac1N}({\rm supp}(f)\cap{\rm supp}(1-\phi)).$$
\end{lemma}
\begin{proof} For a given $n\in\mathbb{Z},$ we have
$$\Big\{\frac{d_Im(I)^{-1}}{|c_I-gc_I|}\|f-f\circ g\|_{L_1(I)}\Big\}_{\substack{I\in\mathbb{D}_{\Gamma}\\ I\cap {\rm supp}(f)\neq\emptyset\\ 3I\cap {\rm supp}(1-\phi)\neq\emptyset\\ d_I>2^{-n}}}\leq c_{\Phi}'\|f\|_{L_1(\mathbb{R}^N)} \Big\{m(I)^{-1}\Big\}_{\substack{I\in\mathbb{D}_{\Gamma}\\ I\cap {\rm supp}(f)\neq\emptyset\\ 3I\cap {\rm supp}(1-\phi)\neq\emptyset}}.$$
The right hand side obviously belongs to $(l_{N,\infty})_0(\mathbb{D}_{\Gamma})$ and, hence, so does the left hand side. Thus,
$${\rm dist}_{l_{N,\infty}(\mathbb{D}_{\Gamma})}\Big(\Big\{\frac{d_Im(I)^{-1}}{|c_I-gc_I|}\|f-f\circ g\|_{L_1(I)}\Big\}_{\substack{I\in\mathbb{D}_{\Gamma}\\ I\cap {\rm supp}(f)\neq\emptyset\\ 3I\cap {\rm supp}(1-\phi)\neq\emptyset}},(l_{N,\infty})_0(\mathbb{D}_{\Gamma})\Big)\leq$$
$$\leq\Big\|\Big\{\frac{d_Im(I)^{-1}}{|c_I-gc_I|}\|f-f\circ g\|_{L_1(I)}\Big\}_{\substack{I\in\mathbb{D}_{\Gamma}\\ I\cap {\rm supp}(f)\neq\emptyset\\ 3I\cap {\rm supp}(1-\phi)\neq\emptyset\\ d_I\leq 2^{-n}}}\Big\|_{l_{N,\infty}(\mathbb{D}_{\Gamma})}.$$
Obviously,
$$|f(x)-f(y)|\leq c_N|x-y|\cdot\|\nabla f\|_{L_{\infty}(\mathbb{R}^N)}.$$
Hence,
$$\|f-f\circ g\|_{L_1(I)}\leq c_N\|\nabla f\|_{L_{\infty}(\mathbb{R}^N)}\times\int_I|x-gx|dx\leq$$
$$\leq c_{\Phi}''m(I)\times |c_I-gc_I|\times \|\nabla f\|_{L_{\infty}(\mathbb{R}^N)}.$$
Thus,
$$\frac{d_Im(I)^{-1}}{|c_I-gc_I|}\|f-f\circ g\|_{L_1(I)}\leq c_{\Phi}'''m(I)^{\frac1N}\|\nabla f\|_{L_{\infty}(\mathbb{R}^N)}$$
Finally,
$${\rm dist}_{l_{N,\infty}(\mathbb{D}_{\Gamma})}\Big(\Big\{\frac{d_Im(I)^{-1}}{|c_I-gc_I|}\|f-f\circ g\|_{L_1(I)}\Big\}_{\substack{I\in\mathbb{D}_{\Gamma}\\ I\cap {\rm supp}(f)\neq\emptyset\\ 3I\cap {\rm supp}(1-\phi)\neq\emptyset}},(l_{N,\infty})_0(\mathbb{D}_{\Gamma})\Big)\leq$$
$$\leq c_{\Phi}\|\nabla f\|_{L_{\infty}(\mathbb{R}^N)}\times\Big\|\Big\{m(I)^{\frac1N}\Big\}_{\substack{I\in\mathbb{D}_{\Gamma}\\ I\cap {\rm supp}(f)\neq\emptyset\\ 3I\cap {\rm supp}(1-\phi)\neq\emptyset\\ d_I\leq 2^{-n}}}\Big\|_{l_{N,\infty}(\mathbb{D}_{\Gamma})}.$$
Passing $n\to\infty$ and using Lemma \ref{first RS extra convergence for sets}, we complete the proof.
\end{proof}

\begin{lemma}\label{first RS final convergence lemma} Let $r>2.$ For every $f\in C^{\infty}_c(\mathbb{R}^N)$ and for every $n\in\mathbb{N},$ we have
$${\rm dist}_{l_{N,\infty}(\mathbb{D}_{\Gamma})}\Big(\Big\{\frac{d_Im(I)^{-\frac1r}}{|c_I-gc_I|}\|f-f\circ g\|_{L_r(I)}\Big\}_{\substack{I\in\mathbb{D}_{\Gamma}\\ I\cap {\rm supp}(f)\neq\emptyset\\ 3I\cap {\rm supp}(1-\phi)\neq\emptyset}},(l_{N,\infty})_0(\mathbb{D}_{\Gamma})\Big)\leq$$
$$\leq c_{\Phi}\|\nabla f\|_{L_{\infty}(\mathbb{R}^N)}\times m^{\frac1N}({\rm supp}(f)\cap{\rm supp}(1-\phi)).$$
\end{lemma}
\begin{proof} We obviously have
$$\|f-f\circ g\|_{L_r(I)}\leq\Big\|f-\frac1{m(I)}\int_If\Big\|_{L_r(I)}+$$
$$+\Big\|f\circ g-\frac1{m(I)}\int_I(f\circ g)\Big\|_{L_r(I)}+m(I)^{\frac1r-1}\|f-f\circ g\|_{L_1(I)}.$$
By Poincare inequality,
$$\|f-f\circ g\|_{L_r(I)}\leq$$
$$\leq c_{\Phi,r}'\Big(m(I)^{\frac1N}\|\nabla f\|_{L_r(I)}+m(I)^{\frac1N}\|\nabla(f\circ g)\|_{L_r(I)}+m(I)^{\frac1r-1}\|f-f\circ g\|_{L_1(I)}\Big).$$
By the triangle inequality,
$${\rm dist}_{l_{N,\infty}(\mathbb{D}_{\Gamma})}\Big(\Big\{\frac{d_Im(I)^{-\frac1r}}{|c_I-gc_I|}\|f-f\circ g\|_{L_r(I)}\Big\}_{\substack{I\in\mathbb{D}_{\Gamma}\\ I\cap {\rm supp}(f)\neq\emptyset\\ 3I\cap {\rm supp}(1-\phi)\neq\emptyset}},(l_{N,\infty})_0(\mathbb{D}_{\Gamma})\Big)\leq$$
$$\leq c_{\Phi,r}''{\rm dist}_{l_{N,\infty}(\mathbb{D}_{\Gamma})}\Big(\Big\{m(I)^{\frac1N-\frac1r}\|\nabla f\|_{L_r(I)}\Big\}_{\substack{I\in\mathbb{D}_{\Gamma}\\ I\cap {\rm supp}(f)\neq\emptyset\\ 3I\cap {\rm supp}(1-\phi)\neq\emptyset}},(l_{N,\infty})_0(\mathbb{D}_{\Gamma})\Big)+$$
$$+c_{\Phi,r}''{\rm dist}_{l_{N,\infty}(\mathbb{D}_{\Gamma})}\Big(\Big\{m(I)^{\frac1N-\frac1r}\|\nabla f\|_{L_r(gI)}\Big\}_{\substack{I\in\mathbb{D}_{\Gamma}\\ I\cap {\rm supp}(f)\neq\emptyset\\ 3I\cap {\rm supp}(1-\phi)\neq\emptyset}},(l_{N,\infty})_0(\mathbb{D}_{\Gamma})\Big)+$$
$$+c_{\Phi,r}''{\rm dist}_{l_{N,\infty}(\mathbb{D}_{\Gamma})}\Big(\Big\{\frac{d_Im(I)^{-1}}{|c_I-gc_I|}\|f-f\circ g\|_{L_1(I)}\Big\}_{\substack{I\in\mathbb{D}_{\Gamma}\\ I\cap {\rm supp}(f)\neq\emptyset\\ 3I\cap {\rm supp}(1-\phi)\neq\emptyset}},(l_{N,\infty})_0(\mathbb{D}_{\Gamma})\Big).$$
The assertion follows from Lemmas \ref{first RS with r convergence lemma} and \ref{first RS without r convergence lemma}.
\end{proof}

\begin{proof}[Proof of Theorem \ref{first RS-type convergence}] The assertion follows by combining Lemmas \ref{first RS convergence lemma} (applied with $f-f\circ g$ instead of $f$) and \ref{first RS final convergence lemma}.
\end{proof}

\section{Second Rochberg-Semmes based estimate}

\begin{theorem}\label{second RS-type theorem} Let $N>1.$ Fix a closed Weyl chamber $\Gamma.$ Suppose that $K:\Gamma\times\Gamma\to\mathbb{C}$ satisfies
$$\bigl|\partial_x^{\alpha}\partial_y^{\beta}K\bigr|(x,y)\leq c_{K,\alpha,\beta}\frac{|x-y|^{-|\alpha|_1-|\beta|_1}}{V(x,|x-y|)},\quad x,y\in\Gamma,\quad\alpha,\beta\in\mathbb{Z}^N_+.$$
Let $V$ be an integral operator with integral kernel $K.$ For every $f\in\dot{W}^{1,N}(\Gamma),$ we have
$$\|[M_f,V]\|_{\mathcal{L}_{N,\infty}(L_2(\Gamma,w(x)dx))}\leq c_{K,\kappa}\|f\|_{\dot{W}^{1,N}(\Gamma)}.$$
\end{theorem}

\begin{theorem}\label{second RS-type convergence} In the setting of Theorem \ref{second RS-type theorem}, let $f\in C^{\infty}_c(\mathbb{R}^N).$ If $0\leq\phi_m\leq 1$ for $m\geq1$ are in $C^{\infty}_c(\mathbb{R}^N)$ and ${\rm supp}(1-\phi_m)$ increases to $\mathbb{R}^N,$ then
$${\rm dist}_{\mathcal{L}_{N,\infty}(L_2(\Gamma,w(x)dx))}\Big(M_{\phi_m}\cdot [M_f,V]\cdot M_{\phi_m}- [M_f,V],(\mathcal{L}_{N,\infty})_0(L_2(\Gamma,w(x)dx))\Big)\to0$$
as $m\to\infty.$
\end{theorem}

\begin{lemma}\label{lem:Weyl-diagonal-kernel-NWO}
Assume that $N>1$ and $2<r<\infty.$ Let $f\in L_r^{\mathrm{loc}}(\mathbb{R}^N).$ Fix a closed Weyl chamber $\Gamma.$ Let $K$ be as in Theorem \ref{first RS-type theorem}. We have
$$(f(x)-f(y))\cdot K(x,y)=\sum_{k,l\in\mathbb{Z}^N}(1+|k|^2+|l|^2)^{-N-1}\sum_{I\in\mathbb{D}_{\Gamma}}e_{k,l,I}^{(1)}(x)f_{k,l,I}^{(1)}(y)-$$
$$-\sum_{k,l\in\mathbb{Z}^N}(1+|k|^2+|l|^2)^{-N-1}\sum_{I\in\mathbb{D}_{\Gamma}}e_{k,l,I}^{(2)}(x)f_{k,l,I}^{(2)}(y)$$
where
\begin{enumerate}[{\rm (i)}]
\item $e_{k,l,I}^{(1)}$ and $e_{k,l,I}^{(2)}$ are supported in $I$ and $f_{k,l,I}^{(1)}$ and $f_{k,l,I}^{(2)}$ are supported in $3I;$
\item for every $k,l\in\mathbb{Z}^N$ and for every $I\in\mathbb{D}_{\Gamma},$ we have
$$\|e_{k,l,I}^{(1)}\|_{L_r(\mathbb{R}^N,w(x)dx)}\leq\Big\|f-\frac1{m(3I)}\int_{3I}f\Big\|_{L_r(3I,w(x)dx)};$$
$$\|f_{k,l,I}^{(2)}\|_{L_r(\mathbb{R}^N,w(x)dx)}=\Big\|f-\frac1{m(3I)}\int_{3I}f\Big\|_{L_r(3I,w(x)dx)};$$
\item for every $k,l\in\mathbb{Z}^N$ and for every $I\in\mathbb{D}_{\Gamma},$ we have
$$\|f_{k,l,I}^{(1)}\|_{L_{\infty}(\mathbb{R}^N)},\|e_{k,l,I}^{(2)}\|_{L_{\infty}(\mathbb{R}^N)}\leq c_{K,\kappa,\Phi};$$
\end{enumerate}
\end{lemma}
\begin{proof} Repeating the argument in Lemma \ref{lem:Weyl-off-diagonal-kernel-NWO} {\it mutatis mutandi}, we obtain
$$(1+|k|^2+|l|^2)^{-N-1}\cdot|{\rm coef}_{k,l,I}|\leq \frac{c_{K,\kappa,\Phi}}{w(I)},\quad k,l\in\mathbb{Z}^N,\quad I\in\mathbb{D}_{\Gamma}.$$
Set
$$e_{k,l,I}^{(1)}(x)=\Big(f(x)-\frac1{m(3I)}\int_{3I}f\Big)\cdot\chi_I(x)e^{i\langle k,\frac{\Phi^{-1}(x-c_I)}{d_I}\rangle},$$
$$f_{k,l,I}^{(1)}(y)=(1+|k|^2+|l|^2)^{N+1}{\rm coef}_{k,l,I}\cdot \chi_{3I}(y)e^{i\langle l,\frac{\Phi^{-1}(y-c_I)}{d_I}\rangle},$$
$$e_{k,l,I}^{(2)}(x)=(1+|k|^2+|l|^2)^{N+1}{\rm coef}_{k,l,I}\cdot\chi_I(x)e^{i\langle k,\frac{\Phi^{-1}(x-c_I)}{d_I}\rangle},$$
$$f_{k,l,I}^{(2)}(y)=\Big(f(y)-\frac1{m(3I)}\int_{3I}f\Big)\cdot \chi_{3I}(y)e^{i\langle l,\frac{\Phi^{-1}(y-c_I)}{d_I}\rangle}.$$
The assertion follows immediately.
\end{proof}

{\color{red} \Huge to be written by Zhijie}

\section{Proof of Theorem \ref{main theorem} \eqref{mta},\eqref{mtc}}

\begin{lemma}\label{heat kernel estimate} For every pair of multi-indices $\alpha,\beta\in\mathbb{Z}^N_+,$ for every $t>0,$ we have
$$|(\partial_x^{\alpha}\partial_y^{\beta}h_t)(x,y)|\leq c_{\kappa,\alpha,\beta}t^{1-\frac{|\alpha|_1+|\beta|_1}{2}}|x-y|^{-2}w(B(x,t^{\frac12}))^{-1}\exp(-\frac{c_{\kappa}}{t}d(x,y)^2).$$
\end{lemma}
\begin{proof} This is (slightly weakened form of) Theorem 3.1 in \cite{MR4042416}.
\end{proof}

\begin{lemma}\label{zhijie kernel estimate} Suppose $s+N>2.$ We have
$$\int_0^{\infty}t^{-\frac{s}{2}}w(B(x,t^{\frac12}))^{-1}\exp(-\frac{c_{\kappa}}{t}d(x,y)^2)dt\leq c_{\kappa,s}w(B(x,d(x,y)))^{-1}\times d(x,y)^{2-s}.$$
\end{lemma}
\begin{proof} By \eqref{doublinginequality},
$$w(B(x,t^{\frac12}))^{-1}\leq c_{\kappa}^{(1)}w(B(x,d(x,y)))^{-1}\times
\begin{cases}
\left(\frac{d(x,y)}{t^{\frac12}}\right)^{\mathbf{N}},
&0<t\leq d(x,y)^2,\\
\left(\frac{d(x,y)}{t^{\frac12}}\right)^N,&t\geq d(x,y)^2.
\end{cases}$$
Therefore,
$$\int_0^{\infty}t^{-\frac{s}{2}}w(B(x,t^{\frac12}))^{-1}\exp(-\frac{c_{\kappa}}{t}d(x,y)^2)dt\leq$$
$$\leq c_{\kappa}^{(1)}w(B(x,d(x,y))^{-1}\times\int_0^{d(x,y)^2}t^{-\frac{s}{2}}\left(\frac{d(x,y)}{t^{\frac12}}\right)^{\mathbf{N}}\exp(-\frac{c_{\kappa}}{t}d(x,y)^2)dt+$$
$$+c_{\kappa}^{(1)}w(B(x,d(x,y)))^{-1}\times\int_{d(x,y)^2}^{\infty}t^{-\frac{s}{2}}\left(\frac{d(x,y)}{t^{\frac12}}\right)^N\exp(-\frac{c_{\kappa}}{t}d(x,y)^2)dt=$$
$$=c_{\kappa}^{(1)}w(B(x,d(x,y)))^{-1}\times d(x,y)^{2-s}\times\int_0^1t^{-\frac{s+\mathbf{N}}{2}}\exp(-\frac{c_{\kappa}}{t})dt+$$
$$+c_{\kappa}^{(1)}w(B(x,d(x,y))/)^{-1}\times d(x,y)^{2-s}\times\int_1^{\infty}t^{-\frac{s+N}{2}}\exp(-\frac{c_{\kappa}}{t})dt.$$
\end{proof}

The following lemma establishes derivative estimates for the integral kernels of Dunkl-Riesz transforms.
\begin{lemma}\label{lem:Dunkl-Riesz-endpoint-derivatives}
Let $N\geq 2.$ For every pair of multi-indices $\alpha,\beta$ and for every $1\leq j\leq N,$ we have
$$|(\partial_x^{\alpha}\partial_y^{\beta}K_{R_{{\rm Dunkl},j}})(x,y)|
\leq c_{\kappa,\alpha,\beta}\frac{d(x,y)^{1-|\alpha|_1-|\beta|_1}}{|x-y|\cdot w(B(x,d(x,y)))},\quad x,y\in\mathbb{R}^N.$$
\end{lemma}
\begin{proof} We have
$$\Delta_{{\rm Dunkl}}^{-\frac12}=\pi^{-\frac12}\int_0^{\infty}e^{-t\Delta_{{\rm Dunkl}}}\frac{dt}{t^{\frac12}}.$$
Hence,
$$K_{\Delta_{{\rm Dunkl}}^{-\frac12}}(x,y)=\pi^{-\frac12}\int_0^{\infty}h_t(x,y)\frac{dt}{t^{\frac12}}.$$
Obviously,
$$K_{R_{{\rm Dunkl},j}}(x,y)=(T_jK_{\Delta_{{\rm Dunkl}}^{-\frac12}})(x,y),$$
where $T_j$ on the right hand side acts in the $x$ variable. Taking into account that
$$(T_jh_t)(x,y)=\frac{(y-x)_j}{2t}h_t(x,y),$$
we conclude that
$$K_{R_{{\rm Dunkl},j}}(x,y)=\frac{(y-x)_j}{2\pi^{\frac12}}\int_0^{\infty}h_t(x,y)\frac{dt}{t^{\frac32}}.$$
By the Leibniz rule,
$$(\partial_x^{\alpha}\partial_y^{\beta}K_{R_{{\rm Dunkl},j}})(x,y)=\frac{(y-x)_j}{2\pi^{\frac12}}\int_0^{\infty}(\partial_x^{\alpha}\partial_y^{\beta}h_t)(x,y)\frac{dt}{t^{\frac32}}-$$
$$-\frac{\alpha_j}{2\pi^{\frac12}}\int_0^{\infty}(\partial_x^{\alpha-e_j}\partial_y^{\beta}h_t)(x,y)\frac{dt}{t^{\frac32}}+\frac{\beta_j}{2\pi^{\frac12}}\int_0^{\infty}(\partial_x^{\alpha}\partial_y^{\beta-e_j}h_t)(x,y)\frac{dt}{t^{\frac32}}.$$
Of course, if $\alpha_j=0$ (respectively, $\beta_j=0$), then the second (respectively, the third) summand on the right hand side does not exist.

Using Lemma \ref{heat kernel estimate}, we estimate
$$|(\partial_x^{\alpha}\partial_y^{\beta}K_{R_{{\rm Dunkl},j}})(x,y)|\leq$$
$$\leq \frac{c_{\kappa,\alpha,\beta}}{2\pi^{\frac12}|x-y|}\int_0^{\infty}t^{-\frac{1+|\alpha|_1+|\beta|_1}{2}}w(B(x,t^{\frac12}))^{-1}\exp(-\frac{c_{\kappa}}{t}d(x,y)^2)dt+$$
$$+\frac{\alpha_jc_{\kappa,\alpha-e_j,\beta}}{2\pi^{\frac12}|x-y|^2}\int_0^{\infty}t^{-\frac{|\alpha|_1+|\beta|_1}{2}}w(B(x,t^{\frac12}))^{-1}\exp(-\frac{c_{\kappa}}{t}d(x,y)^2)dt+$$
$$+\frac{\beta_jc_{\kappa,\alpha,\beta-e_j}}{2\pi^{\frac12}|x-y|^2}\int_0^{\infty}t^{-\frac{|\alpha|_1+|\beta|_1}{2}}w(B(x,t^{\frac12}))^{-1}\exp(-\frac{c_{\kappa}}{t}d(x,y)^2)dt.$$
Applying Lemma \ref{zhijie kernel estimate}, we conclude
$$|(\partial_x^{\alpha}\partial_y^{\beta}K_{R_{{\rm Dunkl},j}})(x,y)|\leq$$
$$\leq \frac{c_{\kappa,\alpha,\beta}}{2\pi^{\frac12}|x-y|}\times c_{\kappa,1+|\alpha|_1+|\beta|_1}w(B(x,d(x,y)))^{-1}\times d(x,y)^{1-|\alpha|_1-|\beta|_1}+$$
$$+\frac{\alpha_jc_{\kappa,\alpha-e_j,\beta}}{2\pi^{\frac12}|x-y|^2}\times c_{\kappa,|\alpha|_1+|\beta|_1}w(B(x,d(x,y)))^{-1}\times d(x,y)^{2-|\alpha|_1-|\beta|_1}+$$
$$+\frac{\beta_jc_{\kappa,\alpha,\beta-e_j}}{2\pi^{\frac12}|x-y|^2}\times c_{\kappa,1+|\alpha|_1+|\beta|_1}w(B(x,d(x,y)))^{-1}\times d(x,y)^{1-|\alpha|_1-|\beta|_1}.$$
The assertion follows immediately.
\end{proof}

\begin{lemma}\label{chamber properties} Let $\Gamma$ be an open chamber and let $1\neq g\in G.$ We have
\begin{enumerate}[{\rm (i)}]
\item $\bar{\Gamma}\cap g\bar{\Gamma}$ is $k$-dimensional chamber of $\Gamma$ for some $1\leq k<d.$ 
\item $g(x)=x$ for $x\in \bar{\Gamma}\cap g\bar{\Gamma}.$
\end{enumerate}
\end{lemma}
\begin{proof} {\color{red} must be a standard material. reference to be found.}
\end{proof}

\begin{lemma}\label{length comparison lemma} Let $N\geq 2,$ $1\leq N_1\leq N$ and $N_2=N-N_1.$ Let $\Omega=\mathbb{R}^{N_1}_+\times\mathbb{R}^{N_2}.$ Let $F$ be a face of $\Omega$ (not necessarily of maximal dimension, possibly of zero dimension). Let $A:\mathbb{R}^N\to\mathbb{R}^N$ be a bijective linear map. Suppose $A(\bar{\Omega})\cap\bar{\Omega}=F,$  $A|_F={\rm id}.$ Suppose also that $A:\mathbb{R}^{N_1}\times\{0\}\to\mathbb{R}^{N_1}\times\{0\},$  $A:\{0\}\times\mathbb{R}^{N_2}\to\{0\}\times\mathbb{R}^{N_2}$ and $A|_{\{0\}\times\mathbb{R}^{N_2}}={\rm id}.$ We have
$$|x-Ay|\geq c_A(|x-y|+{\rm dist}(x,F)),\quad x,y\in\Omega.$$
\end{lemma}
\begin{proof} In what follows, we prove the assertion for the case when $F$ has non-zero dimension. The case ${\rm dim}(F)=0$ (i.e., the case $F=\{0\}$) is similar but simpler.

Denote by $\sigma$ the reflection with respect to $F.$

Since $A,\sigma:\mathbb{R}^{N_1}\times\{0\}\to\mathbb{R}^{N_1}\times\{0\},$  $A,\sigma:\{0\}\times\mathbb{R}^{N_2}\to\{0\}\times\mathbb{R}^{N_2},$ and $A|_{\{0\}\times\mathbb{R}^{N_2}}={\rm id}$ and $\sigma|_{\{0\}\times\mathbb{R}^{N_2}}={\rm id},$ it suffices to prove the assertion for the case $N_1=N,$ $N_2=0.$

Without loss of generality, we may assume that $F=\mathbb{R}^k_+\times\{0\}$ with $1\leq k\leq N.$


Since $A|_{\mathbb{R}^k\times\{0\}}={\rm id},$ it follows that
$$A=
\begin{pmatrix}
1&B\\
0&D
\end{pmatrix}
$$
where $1$ is $k\times k$ matrix, $B$ is $k\times(N-k)$ matrix and $D$ is $(N-k)\times(N-k)$ matrix.

{\bf Step 1:} We claim that $D(\mathbb{R}^{N-k}_+)\cap\mathbb{R}^{N-k}_+=\{0\}.$

Let $w\in\mathbb{R}^{N-k}_+$ be such that $Dw\geq0.$ Set $x=(t,\cdots,t,w)\in\mathbb{R}^N,$ where $t\in\mathbb{R}_+$ is such that $t\geq|Bw|_{\infty}.$ We have $x\geq0$ (this is obvious) and $Ax\geq0$ (because $t(1,\cdots,1)+Bw\geq0$ and $Dw\geq0$). Since $A(\overline{\mathbb{R}^N_+})\cap\overline{\mathbb{R}^N_+}=\mathbb{R}^k_+\times\{0\},$ it follows that $x\in \mathbb{R}^k_+\times\{0\}.$ In other words, $w=0.$

{\bf Step 2:} It follows from Step 1 that $z-Dw\neq0$ for $z,w\in\mathbb{R}^{N-k}_+$ (unless $z=w=0$). Hence,  $|z-Dw|\geq c_D(|z|+|w|)$ for $z,w\in\mathbb{R}^{N-k}_+.$

Set $x=(u,z)$ and $y=(v,w)$ with $u,v\in\mathbb{R}^k_+$ and $z,w\in\mathbb{R}^{N-k}_+.$ We have
$$|x-Ay|=|(u-v-Bw,z-Dw)|\geq |z-Dw|\geq c_D(|z|+|w|).$$
On the other hand, we have
$$|x-Ay|=|(u-v-Bw,z-Dw)|\geq |u-v|-|Bw|\geq |u-v|-\|B\|_{\infty}|w|.$$
Thus,
$$(1+c_D^{-1}\|B\|_{\infty})|x-Ay|\geq |u-v|.$$
Finally,
$$(1+c_D^{-1}\|B\|_{\infty}+c_D^{-1})|x-Ay|\geq |u-v|+|z|+|w|\geq|(u-v,z+w)|=|x-\sigma y|.$$
Obviously,
$$|x-\sigma y|\geq \max\{|x-y|,{\rm dist}(x,F)\}\geq\frac12(|x-y|+{\rm dist}(x,F)).$$
Combining the last two estimates, we complete the proof.
\end{proof}

\begin{proof}[Proof of Theorem \ref{main theorem} \eqref{mtc}] We prove the assertion under {\it a priori} restriction $f\in C^{\infty}_c(\mathbb{R}^N).$ Passing to the case of $f\in\dot{W}^{1,N}(\mathbb{R}^N)$ is standard (see e.g {\color{red} refer to our papers}).

It suffices to estimate, for every chamber $\Gamma$ and for every $g\in G,$ the operator $M_{\chi_{\Gamma}}[R_{{\rm Dunkl},j},M_f]M_{\chi_{\Gamma}}.$ In other words, it suffices to estimate an operator with integral kernel
$$(f(y)-f(x))K_{R_{{\rm Dunkl},j}}(x,y),\quad x\in\Gamma,\quad y\in g\Gamma,$$
on the Hilbert space $L_2(\mathbb{R}^N,w(x)dx).$ In other words, it suffices to estimate an operator with integral kernel
$$(f(gy)-f(x))K_{R_{{\rm Dunkl},j}}(x,gy),\quad x,y\in\Gamma,$$
on the Hilbert space $L_2(\Gamma,w(x)dx).$ 

Let $N_1={\rm dim}({\rm span}(R)),$ $N_2=N-N_1$ and let $\Omega=\mathbb{R}_+^{N_1}\times\mathbb{R}^{N_2}.$ For a concrete chamber $\Gamma,$ let $\Phi:\mathbb{R}^N\to\mathbb{R}^N$ be a linear isomorphism given in Theorem \ref{phi theorem} so that $\Phi(\Gamma)=\Omega.$ Without loss of generality, $\Phi({\rm span}(R))=\mathbb{R}^{N_1}\times\{0\}$ and $\Phi({\rm span}(R)^{\perp})=\{0\}\times\mathbb{R}^{N_2}.$

It suffices to estimate an operator with integral kernel
$$(f(g\Phi^{-1}y)-f(\Phi^{-1}x))K_{R_{{\rm Dunkl},j}}(\Phi^{-1}x,g\Phi^{-1}y),\quad x,y\in\Omega,$$
on the Hilbert space $L_2(\Omega,w(\Phi^{-1}x)dx).$ Using Lemma \ref{s2 trivial lemma}, it suffices to estimate an operator with integral kernel
$$(f(g\Phi^{-1}y)-f(\Phi^{-1}x))K_{R_{{\rm Dunkl},j}}(\Phi^{-1}x,g\Phi^{-1}y),\quad x,y\in\Omega,$$
on the Hilbert space $L_2(\Omega,w_{\theta}(x)dx).$

For $g=1_G,$ we need to estimate an operator with integral kernel
$$(f(\Phi^{-1}y)-f(\Phi^{-1}x))K_{R_{{\rm Dunkl},j}}(\Phi^{-1}x,\Phi^{-1}y),\quad x,y\in\Omega,$$
on the Hilbert space $L_2(\Omega,w_{\theta}(x)dx).$

For $g\in 1_G,$ it suffices to estimate operators with integral kernels
$$(f(g\Phi^{-1}y)-f(g\Phi^{-1}x))K_{R_{{\rm Dunkl},j}}(\Phi^{-1}x,g\Phi^{-1}y),\quad x,y\in\Omega,$$
and, respectively,
$$(f(g\Phi^{-1}x)-f(\Phi^{-1}x))K_{R_{{\rm Dunkl},j}}(\Phi^{-1}x,g\Phi^{-1}y),\quad x,y\in\Omega,$$
on the Hilbert space $L_2(\Omega,w_{\theta}(x)dx).$

Denote by $V_{g,G,j}$ integral operator on the Hilbert space $L_2(\Omega,w_{\theta}(x)dx)$ with integral kernel 
$$K_{g,G,j}(x,y)=K_{R_{{\rm Dunkl},j}}(\Phi^{-1}x,g\Phi^{-1}y),\quad x,y\in\Omega.$$

For $g=1_G,$ it suffices to estimate operator $[M_{f\circ\Phi^{-1}},V_{g,G,j}]$ on the Hilbert space $L_2(\Omega,w_{\theta}(x)dx).$ By Lemmas \ref{lem:Dunkl-Riesz-endpoint-derivatives} and \ref{s2 trivial lemma}, for every pair of multi-indices $\alpha,\beta,$ we have
$$|(\partial_x^{\alpha}\partial_y^{\beta}K_{g,G,j})(x,y)|
\leq c_{\kappa,\alpha,\beta}\frac{|x-y|^{-|\alpha|_1-|\beta|_1}}{w_{\theta}(B(x,|x-y|))},\quad x,y\in\Omega.$$
Applying Theorem \ref{second RS-type theorem} to the operator $V_{g,G,j}$ and to the function $f\circ\Phi^{-1},$ we obtain the required estimate.

For $g\neq 1_G,$ it suffices to estimate operators $[M_{f\circ g\Phi^{-1}},V_{g,G,j}]$ and $M_{f\circ g\Phi^{-1}-f\circ\Phi^{-1}}V_{g,G,j}.$ By Lemmas \ref{lem:Dunkl-Riesz-endpoint-derivatives} and \ref{s2 trivial lemma}, for every pair of multi-indices $\alpha,\beta$ and for every $1\leq j\leq N,$ we have
$$|(\partial_x^{\alpha}\partial_y^{\beta}K_{g,G,j})(x,y)|
\leq c_{\kappa,\alpha,\beta}\frac{|x-y|}{|x-\Phi g\Phi^{-1}y|}\frac{|x-y|^{-|\alpha|_1-|\beta|_1}}{w_{\theta}(B(x,|x-y|))},\quad x,y\in\Omega.$$
By Lemma \ref{chamber properties}, the matrix $\Phi g\Phi^{-1}$ satisfies the assumptions in Lemma \ref{length comparison lemma}. Let $F=\Phi(g\bar{G}\cap G)$ be a face of $\Omega$ (not necessarily of maximal dimension, possibly of zero dimension) as given by Lemma \ref{length comparison lemma}. It follows from Lemma \ref{length comparison lemma} that
$$|(\partial_x^{\alpha}\partial_y^{\beta}K_{g,G,j})(x,y)|
\leq c_{\kappa,\alpha,\beta}\frac{|x-y|}{|x-y|+{\rm dist}(x,F)}\frac{|x-y|^{-|\alpha|_1-|\beta|_1}}{w_{\theta}(B(x,|x-y|))},\quad x,y\in\Omega.$$
Applying Theorem \ref{second RS-type theorem} to the operator $V_{g,G,j}$ and to the function $f\circ g\Phi^{-1},$ we obtain the required estimate for $[M_{f\circ g\Phi^{-1}},V_{g,G,j}].$ Applying Theorem \ref{first RS-type theorem} to the operator $V_{g,G,j}$ and to the function $f\circ g\Phi^{-1}-f\circ\Phi^{-1},$ we obtain the required estimate for $M_{f\circ g\Phi^{-1}-f\circ\Phi^{-1}}V_{g,G,j}.$

This completes the proof.
\end{proof}

\section{Proof of Theorem \ref{main theorem} \eqref{mtb},\eqref{mtd}}\label{lower estimate section}


\begin{lemma}\label{nondegenerate corollary} Let $Q\in\mathcal{D}.$ We have
$$K_j(x,y)\geq \frac{c_k}{w(Q)},\quad x\in 3Q,\quad y\in Q+3s(Q)e_j.$$
\end{lemma}
\begin{proof} {\color{red} TO BE DONE. Follows from Lemma \ref{nondegenerate}.}
\end{proof}

\begin{lemma}\label{improved zhijie lemma} Let real-valued $f\in L_1^{{\rm loc}}(\mathbb{R}^N)$ be such that the commutator $[R_{{\rm Dunkl},j},M_f]$ is bounded on $L_2(\mathbb{R}^N,w(x)dx).$ For every $Q\in\mathcal{D},$ there exist $\xi_Q$ and $\eta_Q$ supported in $7Q$ such that $\|\xi_Q\|_{\infty},\|\eta_Q\|_{\infty}\leq m(Q)^{-\frac12}$ and such that
$$\frac1{w(3Q)}\int_{3Q}|f(y)-f_{3Q}|w(y)dy\leq c_k|\langle[R_{{\rm Dunkl},j},M_f]\xi_Q,\eta_Q\rangle_{L_2(\mathbb{R}^N,w(x)dx)}|.$$
\end{lemma}
\begin{proof} Let $\alpha$ be the median of $f$ on $3Q$ and $\beta$ be the median of $f$ on $Q+3s(Q)e_j.$ Assume for definiteness $\alpha\geq\beta.$ Set
$$A_1=\{x\in 3Q:\ f(x)\geq\alpha\},\quad B_1=\{x\in Q+3s(Q)e_j:\ f(x)\leq\beta\},$$
$$A_2=\{x\in 3Q:\ f(x)\leq\beta\},\quad B_2=\{x\in Q+3s(Q)e_j:\ f(x)\geq\beta\},$$
$$\xi_{Q,s}=w(Q)^{-\frac12}\chi_{A_s},\quad \eta_{Q,s}=w(Q)^{-\frac12}\chi_{B_s},\quad s=1,2.$$

We have
$$\langle[R_{{\rm Dunkl},j},M_f]\xi_{Q,s},\eta_{Q,s}\rangle_{L_2(\mathbb{R}^N,w(x)dx)}=$$
$$=\frac1{w(Q)}\int_{B_s\times A_s}K_j(x,y)(f(y)-f(x))w(x)w(y)dxdy.$$

If $x\in B_1$ and $y\in A_1,$ then
$$f(y)-f(x)=(f(y)-\alpha)+(\alpha-\beta)+(\beta-f(x))=$$
$$=|f(y)-\alpha|+|\alpha-\beta|+|f(x)-\beta|\geq |f(y)-\alpha|+|\alpha-\beta|\geq0.$$
If $x\in B_2$ and $y\in A_2,$ then
$$f(y)-f(x)=(f(y)-\beta)+(\beta-f(x))=$$
$$=-|f(y)-\beta|-|f(x)-\beta|\leq -|f(y)-\beta|\leq 0.$$

Thus,
$$\langle[R_{{\rm Dunkl},j},M_f]\xi_{Q,1},\eta_{Q,1}\rangle_{L_2(\mathbb{R}^N,w(x)dx)}\geq $$
$$\geq \frac1{w(Q)}\int_{B_1\times A_1}K_j(x,y)\big(|f(y)-\alpha|+|\alpha-\beta|\big)w(x)w(y)dxdy\geq $$
$$\geq\frac{c_k}{w^2(Q)}\int_{B_1\times A_1}\big(|f(y)-\alpha|+|\alpha-\beta|\big)w(x)w(y)dxdy=$$
$$=\frac{c_k}{w^2(Q)}\cdot\Big(w(B_1)\cdot\int_{A_1}|f(y)-\alpha|w(y)dy+|\alpha-\beta|w(B_1)w(A_1)\Big)\geq $$
$$=\frac{c_k}{w^2(Q)}\cdot\Big(\frac12w(Q)\cdot\int_{A_1}|f(y)-\alpha|w(y)dy+|\alpha-\beta|\cdot \frac14w^2(Q)\Big)\geq $$
$$\geq\frac{c_k}{4w(Q)}\cdot\Big(\int_{A_1}|f(y)-\alpha|w(y)dy+|\alpha-\beta|w(Q)\Big)$$
and
$$-\langle[R_{{\rm Dunkl},j},M_f]\xi_{Q,2},\eta_{Q,2}\rangle_{L_2(\mathbb{R}^N,w(x)dx)}\geq $$
$$\geq \frac1{w(Q)}\int_{B_2\times A_2}K_j(x,y)|f(y)-\beta|w(x)w(y)dxdy\geq $$
$$\geq\frac{c_k}{w^2(Q)}\int_{B_2\times A_2}|f(y)-\beta|w(x)w(y)dxdy=$$
$$=\frac{c_k}{w^2(Q)}\cdot w(B_2)\cdot\int_{A_2}|f(y)-\beta|w(y)dy\geq $$
$$=\frac{c_k}{2w(Q)}\int_{A_2}|f(y)-\beta|w(y)dy.$$

It remains to note that
$$\int_{3Q}|f(y)-\alpha|w(y)dy=\int_{A_1}|f(y)-\alpha|w(y)dy+$$
$$+\int_{A_2}|f(y)-\alpha|w(y)dy+\int_{3Q\backslash(A_1\cup A_2)}|f(y)-\alpha|w(y)dy\leq$$
$$\leq\int_{A_1}|f(y)-\alpha|w(y)dy+\int_{A_2}|f(y)-\beta|w(y)dy$$
$$+\int_{A_2}|\alpha-\beta|w(y)dy+\int_{3Q\backslash(A_1\cup A_2)}|\alpha-\beta|w(y)dy\leq$$
$$\leq\int_{A_1}|f(y)-\alpha|w(y)dy+\int_{A_2}|f(y)-\beta|w(y)dy+2|\alpha-\beta|w(3Q).$$
Combining the last $3$ equations, we conclude
$$\frac{c_k'}{w(Q)}\int_{3Q}|f(y)-\alpha|w(y)dy\leq\max_{s=1,2}|\langle[R_{{\rm Dunkl},j},M_f]\xi_{Q,s},\eta_{Q,s}\rangle_{L_2(\mathbb{R}^N,w(x)dx)}|.$$

If
$$|\langle[R_{{\rm Dunkl},j},M_f]\xi_{Q,1},\eta_{Q,1}\rangle_{L_2(\mathbb{R}^N,w(x)dx)}|\geq |\langle[R_{{\rm Dunkl},j},M_f]\xi_{Q,2},\eta_{Q,2}\rangle_{L_2(\mathbb{R}^N,w(x)dx)}|,$$
then we set $\xi_Q=\xi_{Q,1}$ and $\eta=\eta_{Q,1}.$ Otherwise, set $\xi_Q=\xi_{Q,2}$ and $\eta_Q=\eta_{Q,2}.$ Thus,
$$\frac{c_k'}{w(Q)}\int_{3Q}|f(y)-\alpha|w(y)dy\leq|\langle[R_{{\rm Dunkl},j},M_f]\xi_Q,\eta_Q\rangle_{L_2(\mathbb{R}^N,w(x)dx)}|.$$
Hence,
$$c_k'|f_{3Q}-\alpha|\leq|\langle[R_{{\rm Dunkl},j},M_f]\xi_Q,\eta_Q\rangle_{L_2(\mathbb{R}^N,w(x)dx)}|$$
and, finally,
$$\frac{c_k'}{w(Q)}\int_{3Q}|f(y)-f_{3Q}|w(y)dy\leq 2|\langle[R_{{\rm Dunkl},j},M_f]\xi_Q,\eta_Q\rangle_{L_2(\mathbb{R}^N,w(x)dx)}|.$$
\end{proof}

\begin{lemma}\label{standard} Suppose that $1<p<\infty$ and $f\in L_1^{{\rm loc}}(\mathbb{R}^N).$ If $[R_{{\rm Dunkl},j},M_f]\in\mathcal{L}_p(L_2(\mathbb{R}^N,w(x)dx))$ for some $j\in\{1,2,\cdots,N\},$ then
$$\Big\|\Big\{\frac1{w(3Q)}\int_{3Q}|f(x)-f_{3Q}|w(x)dx\Big\}_{Q\in\mathcal{D}}\Big\|_{l_p(\mathcal{D})}\leq$$
$$\leq c_k\|[R_{{\rm Dunkl},j},M_f]\|_{\mathcal{L}_p(L_2(\mathbb{R}^N,w(x)dx))}.$$
If $[R_{{\rm Dunkl},j},M_f]\in\mathcal{L}_{N,\infty}(L_2(\mathbb{R}^N,w(x)dx))$ for some $j\in\{1,2,\cdots,N\},$ then
$$\Big\|\Big\{\frac1{w(3Q)}\int_{3Q}|f(x)-f_{3Q}|w(x)dx\Big\}_{Q\in\mathcal{D}}\Big\|_{l_{N,\infty}(\mathcal{D})}\leq$$
$$\leq c_k\|[R_{{\rm Dunkl},j},M_f]\|_{\mathcal{L}_{N,\infty}(L_2(\mathbb{R}^N,w(x)dx))}.$$
\end{lemma}
\begin{proof} We prove only the first assertion. The proof of the second one follows {\it mutatis mutandi}.

Let $\{\xi_Q\}_{Q\in\mathcal{D}}$ and $\{\eta_Q\}_{Q\in\mathcal{D}}$ be functions constructed in Lemma \ref{improved zhijie lemma}. By Lemma \ref{NWO lower estimate},
$$\Big\{\langle [R_{{\rm Dunkl},j},M_f]\xi_Q,\eta_Q\rangle_{L_2(\mathbb{R}^N,w(x)dx)}\Big\}_{Q\in\mathcal{D}}\in l_p.$$
The first assertion follows now from Lemma \ref{improved zhijie lemma}.
\end{proof}

\begin{lemma}\label{for reverse holder first lemma} Let $Q$ be a cube in $\mathbb{R}^N$ (with sides parallel to coordinate hyperplanes). We have
$$\sup_{x\in Q}w(x)\leq \frac{c_k}{m(Q)}\int_Qw(x)dx.$$
\end{lemma}
\begin{proof} If $x=c_Q+y,$ $y\in [-\frac12s(Q),\frac12s(Q)]^d,$ then
$$|\langle x,v\rangle|\leq |\langle c_Q,v\rangle|+r|\langle y,v\rangle|\leq |\langle c_Q,v\rangle|+\frac{\sqrt{d}}{2}s(Q).$$
Thus,
$$w(x)=\prod_{v\in R}|\langle x,v\rangle|^{k(v)}\leq\prod_{v\in R}(|\langle c_Q,v\rangle|+\frac{\sqrt{d}}{2}s(Q))^{k(v)}.$$
In particular,
$$LHS\leq \prod_{v\in R}(|\langle c_Q,v\rangle|+\frac{\sqrt{d}}{2}s(Q))^{k(v)}.$$

Fix chamber $\Gamma$ such that $c_Q\in\Gamma.$ We have
$$\int_Qw(x)dx\geq\int_{c_Q+(B(0,\frac12s(Q))\cap\Gamma)}w(x)dx=\int_{B(0,\frac12s(Q))\cap\Gamma}w(c_Q+y)dy.$$
If $y\in\Gamma,$ then
$$|\langle c_Q+y,v\rangle|=|\langle c_Q,v\rangle|+|\langle y,v\rangle|,\quad v\in R.$$
Thus,
$$w(c_Q+y)=\prod_{v\in R}(|\langle c_Q,v\rangle|+|\langle y,v\rangle|)^{k(v)}.$$
Hence,
$$\int_Qw(x)dx\geq\int_{B(0,\frac12s(Q))\cap\Gamma}\prod_{v\in R}(|\langle c_Q,v\rangle|+|\langle y,v\rangle|)^{k(v)}dy.$$
Fix a closed cone $\Gamma'\subset\Gamma$ such that (i) interior of $\Gamma'$ is non-empty and (ii) $\partial\Gamma\cap\partial\Gamma'=\{0\}.$ We have
$$|\langle y,v\rangle|\geq c_{\Gamma'}|y|,\quad y\in\Gamma',\quad v\in R.$$
We, therefore, have
$$\int_Qw(x)dx\geq\int_{B(0,\frac12s(Q))\cap\Gamma'}\prod_{v\in R}(|\langle c_Q,v\rangle|+c_{\Gamma'}|y|)^{k(v)}dy\geq$$
$$\geq \int_{\substack{y\in\Gamma'\\ \frac14s(Q)\leq|y|\leq \frac14s(Q)}}\prod_{v\in R}(|\langle c_Q,v\rangle|+\frac12c_{\Gamma'}s(Q))^{k(v)}dy=$$
$$=s(Q)^N\prod_{v\in R}(|\langle c_Q,v\rangle|+\frac14c_{\Gamma'}s(Q))^{k(v)}\times (2^{-N}-4^{-N})\int_{\substack{y\in\Gamma'\\ |y|\leq1}}dy.$$
Combining this with the preceding paragraph, we complete the proof.
\end{proof}

\begin{lemma}\label{for reverse holder second lemma} Suppose $1<t<1+\frac1{|R|\cdot\max_{v\in R}k(v)}.$ Let $Q$ be a cube in $\mathbb{R}^N$ (with sides parallel to coordinate hyperplanes). We have
$$\Big(\frac{\int_Qw^{1-t}(x)dx}{\int_Qw(x)dx}\Big)^{\frac1t}\leq c_{t,k} \frac{m(Q)}{\int_Qw(x)dx}.$$
\end{lemma}
\begin{proof} If $v\in R$ is such that
$$\sqrt{d}s(Q)\leq |\langle c_Q,v\rangle|,$$
then
$$|\langle x,v\rangle|\geq |\langle c_Q,v\rangle|-|\langle x-c_Q,v\rangle|\geq |\langle c_Q,v\rangle|-\frac{\sqrt{d}}{2}s(Q)\geq \frac12|\langle c_Q,v\rangle|,\quad x\in Q.$$
In this case,
$$\int_Q|\langle x,v\rangle|^{(1-t)|R|k(v)}dx\leq 2^{(t-1)|R|k(v)}|\langle c_Q,v\rangle|^{(1-t)|R|k(v)}\cdot\int_Qdx=$$
$$=2^{(t-1)|R|k(v)}|\langle x_0,v\rangle|^{(1-t)|R|k(v)}\cdot s(Q)^N.$$

If $v\in R$ is such that
$$|\langle c_Q,v\rangle|\leq \sqrt{d}s(Q),$$
then let $y_0=\langle c_Q,v\rangle$ and $z_0=c_Q-y_0v.$ Let us introduce new variables $y=\langle x,v\rangle$ and $z=x-yv$ (it varies in the plane orthogonal to $v$). We write
$$\int_Q|\langle x,v\rangle|^{(1-t)|R|k(v)}dx\leq\int_{|x-c_Q|\leq\frac{\sqrt{d}}{2}s(Q)}|\langle x,v\rangle|^{(1-t)|R|k(v)}dx\leq$$
$$\leq \int_{|y-y_0|,|z-z_0|\leq \frac{\sqrt{d}}{2}r}|y|^{(1-t)|R|k(v)}dydz=$$
$$=(\frac{\sqrt{d}}{2}s(Q))^{N-1}{\rm Vol}(\mathbb{B}^{N-1})\cdot\int_{|y-y_0|\leq \frac{\sqrt{d}}{2}s(Q)}|y|^{(1-t)|R|k(v)}dy\leq$$
$$=(\frac{\sqrt{d}}{2}s(Q))^{N-1}{\rm Vol}(\mathbb{B}^{N-1})\cdot\int_{|y|\leq \frac{3\sqrt{d}}{2}s(Q)}|y|^{(1-t)|R|k(v)}dy=$$
$$=s(Q)^{N+(1-t)|R|k(v)}\cdot (\frac{\sqrt{d}}{2})^{N-1}{\rm Vol}(\mathbb{B}^{N-1})\cdot\int_{|y|\leq \frac{3\sqrt{d}}{2}}|y|^{(1-t)|R|k(v)}dy.$$

Combining the estimates in the preceding paragraphs, we write
$$\int_Q|\langle x,v\rangle|^{(1-t)|R|k(v)}dx\leq c_1s(Q)^N\big(\max\{|\langle c_Q,v\rangle|,r\}\big)^{(1-t)|R|k(v)}.$$
By H\"older inequality,
$$\int_Qw^{1-t}(x)dx\leq \prod_{v\in R}\Big(\int_Q|\langle x,v\rangle|^{(1-t)|R|k(v)}dx\Big)^{\frac1{|R|}}\leq$$
$$\leq c_2s(Q)^N\prod_{v\in R}\big(\max\{|\langle c_Q,v\rangle|,s(Q)\}\big)^{(1-t)k(v)}.$$
The assertion follows easily.
\end{proof}

\begin{lemma}\label{hytonen verification lemma} Measures $dx$ and $w(x)dx$ (and distance $(x,y)\to |x-y|_{\infty}$) satisfy the conditions in Proposition 1.2 in \cite{Hytonen-new}.
\end{lemma}
\begin{proof} That both measures are doubling is obvious.

Let $Q$ be a cube in $\mathbb{R}^N$ (with sides parallel to coordinate hyperplanes). By Lemmas \ref{for reverse holder first lemma} and \ref{for reverse holder second lemma},
$$\Big(\frac1{m(Q)}\int_Qw^t(x)dx\Big)^{\frac1t}\leq\frac{c_k}{m(Q)}\int_Qw(x)dx,\quad t>0.$$
$$\Big(\frac{\int_Qw^{1-t}(x)dx}{\int_Qw(x)dx}\Big)^{\frac1t}\leq c_G \frac{m(Q)}{\int_Qw(x)dx},\quad 1<t<1+\frac1{|R|\cdot\max_{v\in R}k(v)}.$$
The assertion follows now from Proposition 3.6 in \cite{Hytonen-new}.
\end{proof}

\begin{proof}[Proof of Theorem \ref{main theorem} \eqref{mtb}] The left hand side in Lemma \ref{standard} is the oscillation semi-norm of the function $f$ with respect to measure $w(x)dx$ and distance $(x,y)\to |x-y|_{\infty}.$ So, the oscillation semi-norm of the function $f$ with respect to measure $w(x)dx$ and distance $(x,y)\to |x-y|_{\infty}$ is controlled by the $\|\cdot\|_{\mathcal{L}_p(L_2(\mathbb{R}^N,w(x)dx))}$-norm of the commutator $[R_{{\rm Dunkl},j},M_f].$

Lemma \ref{hytonen verification lemma} asserts that measures $dx$ and $w(x)dx$ (and distance $(x,y)\to |x-y|_{\infty}$) satisfy the conditions in Proposition 1.2 in \cite{Hytonen-new}. Appealing to Proposition 1.2 in \cite{Hytonen-new}, we conclude that the oscillation semi-norm of the function $f$ with respect to measure $dx$ and distance $(x,y)\to |x-y|_{\infty}$ is controlled by the one with respect to measure $w(x)dx$ and distance $(x,y)\to |x-y|_{\infty}.$ By the preceding paragraph, the oscillation semi-norm of the function $f$ with respect to measure $dx$ and distance $(x,y)\to |x-y|_{\infty}$ is controlled by the $\|\cdot\|_{\mathcal{L}_p(L_2(\mathbb{R}^N,w(x)dx))}$-norm of the commutator $[R_{{\rm Dunkl},j},M_f].$

That oscillation semi-norm of the function $f$ with respect to measure $dx$ and distance $(x,y)\to |x-y|_{\infty}$ is equivalent to the Sobolev semi-norm of the function $f$ is a standard fact. The proof (in much more general setting) is available in Proposition 1.29 in \cite{HKarxiv}. The assertion follows now from the preceding paragraph. 
\end{proof}

\begin{proof}[Proof of Theorem \ref{main theorem} \eqref{mtd}] The argument follows the one in the proof of Theorem \ref{main theorem} \eqref{mtb} {\it mutatis mutandi}.
\end{proof}


\section{Spectral asymptotic formula}
\setcounter{equation}{0}

\begin{definition}
For \(g\in G\) and \(j=1,\ldots,N\), set
\begin{align*}
d_{g,j}:=g^{-1}e_j,
\qquad
f_g(x):=f(gx).
\end{align*}
For \(\alpha\in R_+\), define \(S_\alpha:\mathbb C^G\to\mathbb C^G\) by
\begin{align*}
(S_\alpha z)_g:=z_{g\sigma_\alpha},
\qquad g\in G.
\end{align*}
For \(f\in\dot W^{1,N}(\mathbb R^N)\), \(x\in\Gamma_+\), and
\(\theta\in\mathbb S^{N-1}\), define
\begin{align*}
\mathfrak c_{f,j}(x,\theta)
&:=
\operatorname{diag}_{g\in G}
\bigg(
\langle d_{g,j},\nabla f_g(x)\rangle
-
\langle d_{g,j},\theta\rangle
\langle\theta,\nabla f_g(x)\rangle
\bigg)\\
&\quad+
\sum_{\alpha\in R_+}
k(\alpha)
\operatorname{diag}_{g\in G}
\bigg(
\frac{\langle\alpha,d_{g,j}\rangle}
{\langle\alpha,x\rangle}
\big(
f_g(x)-f_{g\sigma_\alpha}(x)
\big)
\bigg)
S_\alpha.
\end{align*}
For smooth \(f\), the quotients in the second term are understood by their smooth extensions across the corresponding reflection hyperplanes. For \(f\in\dot W^{1,N}(\mathbb R^N)\), \(\mathfrak c_{f,j}\) is understood almost everywhere on \(\Gamma\times\mathbb S^{N-1}\).
\end{definition}

\begin{lemma}\label{lem:Dunkl-T-commutator-exact}
Let \(f\in W_{\mathrm{loc}}^{1,1}(\mathbb R^N)\). For every \(u\in C_c^\infty(\mathbb R^N)\),
\begin{align}
[T_{e_j},M_f]u(x)
&=
\partial_j f(x)u(x)
+
\sum_{\alpha\in R_+}
k(\alpha)\langle\alpha,e_j\rangle
\frac{f(x)-f(\sigma_\alpha x)}
{\langle\alpha,x\rangle}
u(\sigma_\alpha x)
\label{eq:Dunkl-T-commutator-exact}
\end{align}
for almost every \(x\in\mathbb R^N\). If \(N>1\) and \(f\in\dot W^{1,N}(\mathbb R^N)\), then, for every \(\alpha\in R_+\),
\begin{align}
\left\|
\frac{f-f\circ\sigma_\alpha}
{\langle\alpha,\cdot\rangle}
\right\|_{L^N(\mathbb R^N,dx)}
\lesssim_{N,G}
\|\nabla f\|_{L^N(\mathbb R^N,dx)}.
\label{eq:Dunkl-reflection-difference-quotient}
\end{align}
\end{lemma}

\begin{proof}
Pairing \(\alpha\) and \(-\alpha\) in the definition of \(T_{e_j}\), and using \(k(-\alpha)=k(\alpha)\) and \(\sigma_{-\alpha}=\sigma_\alpha\), we have
\begin{align*}
T_{e_j}u(x)
=
\partial_j u(x)
+
\sum_{\alpha\in R_+}
k(\alpha)\langle\alpha,e_j\rangle
\frac{u(x)-u(\sigma_\alpha x)}
{\langle\alpha,x\rangle}.
\end{align*}
Therefore,
\begin{align*}
[T_{e_j},M_f]u(x)
&=
\partial_jf(x)u(x)
+
\sum_{\alpha\in R_+}
k(\alpha)\langle\alpha,e_j\rangle
\frac{f(x)-f(\sigma_\alpha x)}
{\langle\alpha,x\rangle}
u(\sigma_\alpha x),
\end{align*}
which proves \eqref{eq:Dunkl-T-commutator-exact}.

Since \(R\) is normalized,
\begin{align*}
x-\sigma_\alpha x
=
\langle\alpha,x\rangle\alpha,
\qquad
\|x-\sigma_\alpha x\|
=
\sqrt{2}\,|\langle\alpha,x\rangle|.
\end{align*}
Hence, for almost every \(x\notin\alpha^\perp\),
\begin{align*}
\frac{|f(x)-f(\sigma_\alpha x)|}
{|\langle\alpha,x\rangle|}
=
\sqrt{2}\,
\mathcal G_{I,\sigma_\alpha}f(x).
\end{align*}
By Lemma \ref{lem:Weyl-difference-quotient} \eqref{eq:Weyl-difference-quotient},
\begin{align*}
\left\|
\frac{f-f\circ\sigma_\alpha}
{\langle\alpha,\cdot\rangle}
\right\|_{L^N(\mathbb R^N,dx)}
\lesssim_{N,G}
\|\nabla f\|_{L^N(\mathbb R^N,dx)}.
\end{align*}
This completes the proof of Lemma \ref{lem:Dunkl-T-commutator-exact}.
\end{proof}
\begin{lemma}\label{lem:Dunkl-chamber-folding-symbol}
Let \(\Gamma_+\) be an open Weyl chamber, let \(\Gamma\) be its closure, and let \(R_+\) be the corresponding positive subsystem. Set
\begin{align*}
\rho(x):=\varpi(x)^{1/2},
\end{align*}
and define
\begin{align*}
U:L^2(\mathbb R^N,dw)&\longrightarrow L^2(\Gamma,dw;\mathbb C^G),
&
(Uu)_g(x)&:=u(gx),\\
V:L^2(\Gamma,dw;\mathbb C^G)&\longrightarrow L^2(\Gamma,dx;\mathbb C^G),
&
(Vh)(x)&:=\rho(x)h(x).
\end{align*}
Then \(U\) and \(V\) are unitary. If
\begin{align*}
\widetilde T_j
:=
VUT_{e_j}U^{-1}V^{-1},
\qquad
\widetilde P
:=
VU(-\Delta_{{\rm Dunkl}})U^{-1}V^{-1},
\end{align*}
then, on \(C_c^\infty(\Gamma_+;\mathbb C^G)\),
\begin{align}
(\widetilde T_jh)_g
&=
\partial_{d_{g,j}}h_g
-
\sum_{\alpha\in R_+}
k(\alpha)
\frac{\langle\alpha,d_{g,j}\rangle}
{\langle\alpha,x\rangle}
h_{g\sigma_\alpha},
\label{eq:Dunkl-folded-T}\\
\widetilde P
&=
-\Delta I_{\mathbb C^G}
+
|\nabla\log\rho|^2I_{\mathbb C^G}
-
2\sum_{\alpha\in R_+}
k(\alpha)
\frac{S_\alpha}{\langle\alpha,x\rangle^2}.
\label{eq:Dunkl-folded-Laplacian}
\end{align}

Assume in addition that \(N>1\), and let
\begin{align*}
\widetilde Q
:=
VU(-\Delta_{{\rm Dunkl}})^{-1/2}U^{-1}V^{-1},
\end{align*}
where \((-\Delta_{{\rm Dunkl}})^{-1/2}\) is defined by the spectral calculus.
Then, for every \(\chi\in C_c^\infty(\Gamma_+)\), the operator
\(M_\chi\widetilde Q M_\chi\), initially defined on its natural
domain, admits a bounded extension to
\(L^2(\Gamma,dx;\mathbb C^G)\), and this extension is a classical
matrix-valued pseudodifferential operator of order \(-1\). Its
principal symbol is
\begin{align*}
\sigma_{-1}
\big(
M_\chi\widetilde Q M_\chi
\big)(x,\xi)
=
\chi(x)^2|\xi|^{-1}I_{\mathbb C^G}.
\end{align*}
Consequently, if \(x\in\Gamma_+\), \(\xi\neq0\), and
\(\chi\in C_c^\infty(\Gamma_+)\) is identically \(1\) in a
neighborhood of \(x\), then
\begin{align}
\sigma_{-1}
\big(
M_\chi\widetilde Q M_\chi
\big)(x,\xi)
=
|\xi|^{-1}I_{\mathbb C^G}.
\label{eq:Dunkl-folded-Q-principal-symbol}
\end{align}
\end{lemma}
\begin{proof}
The Weyl chambers \(g\Gamma_+\), \(g\in G\), are pairwise disjoint
and cover \(\mathbb R^N\) up to the reflection hyperplanes; see
\cite[Theorem~2.12]{HeckmanCoxeterGroups}. Since \(dw\) is
\(G\)-invariant,
\begin{align*}
\|Uu\|_{L^2(\Gamma,dw;\mathbb C^G)}^2
&=
\sum_{g\in G}
\int_\Gamma
|u(gx)|^2\,dw(x)\\
&=
\int_{\mathbb R^N}|u(x)|^2\,dw(x).
\end{align*}
Conversely, for
\(h=(h_g)_{g\in G}\in L^2(\Gamma,dw;\mathbb C^G)\), define
\(u(gx):=h_g(x)\) for \(x\in\Gamma_+\). This determines \(u\) almost
everywhere on \(\mathbb R^N\) and gives \(Uu=h\). Hence \(U\) is
unitary. Since \(dw=\rho^2dx\) and \(\rho>0\) on \(\Gamma_+\),
\(V\) is unitary as well.

Pairing the roots \(\alpha\) and \(-\alpha\) in the definition of
\(T_{e_j}\), changing variables in the root system by \(g\), and using
the \(G\)-invariance of \(k\), we obtain, for
\(h\in C_c^\infty(\Gamma_+;\mathbb C^G)\),
\begin{align*}
(UT_{e_j}U^{-1}h)_g
=
\partial_{d_{g,j}}h_g
+
\sum_{\alpha\in R_+}
k(\alpha)
\frac{\langle\alpha,d_{g,j}\rangle}
{\langle\alpha,x\rangle}
\bigl(h_g-h_{g\sigma_\alpha}\bigr).
\end{align*}
On \(\Gamma_+\),
\begin{align*}
\nabla\log\rho(x)
=
\sum_{\alpha\in R_+}
k(\alpha)\frac{\alpha}{\langle\alpha,x\rangle},
\end{align*}
and therefore
\begin{align*}
\rho\,\partial_{d_{g,j}}\rho^{-1}
=
\partial_{d_{g,j}}
-
\sum_{\alpha\in R_+}
k(\alpha)
\frac{\langle\alpha,d_{g,j}\rangle}
{\langle\alpha,x\rangle}.
\end{align*}
Conjugation by \(V\) cancels the diagonal terms and gives
\eqref{eq:Dunkl-folded-T}.

Similarly, pairing \(\alpha\) and \(-\alpha\) in the formula for the
Dunkl Laplacian and using \(\|\alpha\|^2=2\), we have
\begin{align*}
\bigl(
U(-\Delta_{{\rm Dunkl}})U^{-1}h
\bigr)_g
&=
-\Delta h_g
-
2\sum_{\alpha\in R_+}
k(\alpha)
\frac{\partial_\alpha h_g}
{\langle\alpha,x\rangle}\\
&\quad+
2\sum_{\alpha\in R_+}
k(\alpha)
\frac{h_g-h_{g\sigma_\alpha}}
{\langle\alpha,x\rangle^2}.
\end{align*}
Writing \(\phi=\log\rho\),
\begin{align*}
\rho
\bigl(
-\Delta-2\nabla\phi\cdot\nabla
\bigr)
\rho^{-1}
=
-\Delta+|\nabla\phi|^2+\Delta\phi,
\end{align*}
while
\begin{align*}
\Delta\phi
=
-\sum_{\alpha\in R_+}
k(\alpha)
\frac{\|\alpha\|^2}{\langle\alpha,x\rangle^2}
=
-2\sum_{\alpha\in R_+}
\frac{k(\alpha)}{\langle\alpha,x\rangle^2}.
\end{align*}
The latter term cancels the diagonal part of
\[
2\sum_{\alpha\in R_+}
k(\alpha)
\frac{I_{\mathbb C^G}-S_\alpha}
{\langle\alpha,x\rangle^2},
\]
and hence
\[
\widetilde P
=
-\Delta I_{\mathbb C^G}
+
|\nabla\log\rho|^2I_{\mathbb C^G}
-
2\sum_{\alpha\in R_+}
k(\alpha)
\frac{S_\alpha}{\langle\alpha,x\rangle^2},
\]
which proves \eqref{eq:Dunkl-folded-Laplacian}.

Assume now that \(N>1\), fix
\(\chi\in C_c^\infty(\Gamma_+)\), and set
\[
K:=\operatorname{supp}\chi.
\]
On \(\Gamma_+\), write
\[
\widetilde P=-\Delta I_{\mathbb C^G}+W(x),
\qquad
W(x)
=
|\nabla\log\rho(x)|^2I_{\mathbb C^G}
-
2\sum_{\alpha\in R_+}
k(\alpha)
\frac{S_\alpha}{\langle\alpha,x\rangle^2}.
\]
Thus \(W\) is a smooth Hermitian matrix-valued function on
\(\Gamma_+\).

Choose open sets
\[
K\Subset\Omega_0\Subset\Omega_1\Subset\Gamma_+
\]
and \(\eta\in C_c^\infty(\Gamma_+)\) such that
\[
0\leq\eta\leq1,
\qquad
\eta=1
\quad\text{on a neighborhood of }\overline{\Omega_1}.
\]
Choose a flat torus \(M\) containing
\(\operatorname{supp}\eta\) in one Euclidean coordinate chart.
Since \(\eta\) is compactly supported in that chart, the differential
operator \(M_\eta\widetilde P M_\eta\) extends smoothly by zero to
\(M\). Define on \(L^2(M;\mathbb C^G)\)
\begin{align}
P_0
:=
M_\eta\widetilde P M_\eta
+
M_{1-\eta}(-\Delta_M+1)M_{1-\eta}.
\label{eq:Dunkl-local-elliptic-extension}
\end{align}
It is a nonnegative self-adjoint elliptic differential operator of
order \(2\). Indeed,
\begin{align*}
\langle P_0u,u\rangle
&=
\langle
\widetilde P(\eta u),\eta u
\rangle
+
\|\nabla((1-\eta)u)\|_2^2
+
\|(1-\eta)u\|_2^2
\geq0,
\end{align*}
and its principal symbol is
\begin{align*}
p_{0,2}(x,\xi)
=
\bigl(
\eta(x)^2+(1-\eta(x))^2
\bigr)
|\xi|^2I_{\mathbb C^G},
\end{align*}
which is elliptic since
\[
\eta^2+(1-\eta)^2\geq\frac12.
\]
Moreover,
\begin{align}
P_0=\widetilde P
\quad\text{on }\Omega_1.
\label{eq:Dunkl-P-P0-local-equality}
\end{align}

Let
\[
H_t:=e^{-t\widetilde P},
\qquad
H_t^0:=e^{-tP_0}.
\]
We first consider
\begin{align*}
B_0
:=
\frac1{\sqrt\pi}
\int_0^1
t^{-1/2}H_t^0\,dt.
\end{align*}
Let \(\Pi_0\) be the orthogonal projection onto \(\ker P_0\).
Since \(P_0\) is elliptic on the closed manifold \(M\),
\(\Pi_0\) is smoothing. Applying the complex powers calculus
\cite[Theorem~2 and p.~290]{MR0237943} to
\(P_0+\Pi_0\), we obtain
\begin{align*}
(P_0+\Pi_0)^{-1/2}
\in
\Psi_{\mathrm{cl}}^{-1}(M;\mathbb C^G)
\end{align*}
with principal symbol \(p_{0,2}^{-1/2}\).

Set
\[
F(\lambda)
:=
\frac1{\sqrt\pi}
\int_0^1 t^{-1/2}e^{-t\lambda}\,dt,
\qquad
\lambda\geq0.
\]
Then \(B_0=F(P_0)\), and for \(\lambda>0\),
\begin{align*}
F(\lambda)-\lambda^{-1/2}
=
-\frac1{\sqrt\pi}
\int_1^\infty
t^{-1/2}e^{-t\lambda}\,dt.
\end{align*}
Hence
\[
F(\lambda)-\lambda^{-1/2}
=
O(\lambda^{-m})
\qquad
(\lambda\to\infty)
\]
for every \(m>0\). On \(\ker P_0\) the difference between \(B_0\)
and \((P_0+\Pi_0)^{-1/2}\) is finite rank. It follows from elliptic
spectral regularity that
\begin{align*}
B_0-(P_0+\Pi_0)^{-1/2}
\in\Psi^{-\infty}(M;\mathbb C^G).
\end{align*}
Consequently,
\begin{align}
B_0\in\Psi_{\mathrm{cl}}^{-1}(M;\mathbb C^G),
\qquad
\sigma_{-1}(B_0)(x,\xi)
=
p_{0,2}(x,\xi)^{-1/2}.
\label{eq:Dunkl-B0-symbol}
\end{align}
In particular, by \eqref{eq:Dunkl-P-P0-local-equality},
\begin{align}
\sigma_{-1}(B_0)(x,\xi)
=
|\xi|^{-1}I_{\mathbb C^G},
\qquad
x\in\Omega_1.
\label{eq:Dunkl-B0-local-symbol}
\end{align}

We next compare \(H_t\) and \(H_t^0\) in the interior. Choose
\(\zeta\in C_c^\infty(\Omega_1)\) such that
\[
\zeta=1
\quad\text{on a neighborhood of }K.
\]
Since \(P_0=\widetilde P\) on \(\Omega_1\), Duhamel's formula gives
\begin{align}
M_\chi(H_t-H_t^0)M_\chi
=
-\int_0^t
M_\chi H_{t-s}
[\widetilde P,M_\zeta]
H_s^0M_\chi\,ds.
\label{eq:Dunkl-local-heat-Duhamel}
\end{align}
Indeed, for data supported in \(K\),
\[
v(t):=M_\zeta H_t^0M_\chi u
\]
satisfies
\[
(\partial_t+\widetilde P)v
=
[\widetilde P,M_\zeta]H_t^0M_\chi u,
\qquad
v(0)=M_\chi u,
\]
which yields \eqref{eq:Dunkl-local-heat-Duhamel}.

Let \(h_t(x,y)\) be the Dunkl heat kernel. The integral kernel of
\(H_t\), with respect to Lebesgue measure on \(\Gamma\), is the
matrix
\begin{align}
\widetilde h_t(x,y)_{g,\ell}
=
\rho(x)\rho(y)\,
h_t(gx,\ell y),
\qquad g,\ell\in G.
\label{eq:Dunkl-folded-heat-kernel}
\end{align}
For \(x,y\in\Gamma\),
\[
d(gx,\ell y)=d(x,y)=|x-y|.
\]
Together with the mixed derivative estimate for the Dunkl heat
kernel \cite[Theorem~3.1]{MR4042416}, this implies that, whenever
\(E,F\Subset\Gamma_+\) and
\(\operatorname{dist}(E,F)>0\), for every pair of multi-indices
\(\alpha,\beta\) and every \(m>0\),
\begin{align}
\sup_{\substack{x\in E\\y\in F}}
\left\|
\partial_x^\alpha
\partial_y^\beta
\widetilde h_t(x,y)
\right\|
=
O(t^m),
\qquad
t\downarrow0.
\label{eq:Dunkl-folded-heat-off-diagonal}
\end{align}
The heat kernel of the smooth elliptic operator \(P_0\) satisfies the
same off-diagonal estimate on \(M\).

Since
\[
\operatorname{supp}[\widetilde P,M_\zeta]
\cap K
=
\varnothing,
\]
the two heat kernels occurring in
\eqref{eq:Dunkl-local-heat-Duhamel} are evaluated between compact
sets separated by a positive distance. The standard Gaussian
bounds underlying
\eqref{eq:Dunkl-folded-heat-off-diagonal} therefore give, for every
multi-indices \(\alpha,\beta\) and every \(m>0\),
\begin{align}
\sup_{x,y}
\left\|
\partial_x^\alpha\partial_y^\beta
K_{M_\chi(H_t-H_t^0)M_\chi}(x,y)
\right\|
=
O(t^m),
\qquad
t\downarrow0.
\label{eq:Dunkl-local-heat-flat}
\end{align}
Hence
\begin{align}
S_0
:=
\frac1{\sqrt\pi}
\int_0^1
t^{-1/2}
M_\chi(H_t-H_t^0)M_\chi\,dt
\in
\Psi^{-\infty}.
\label{eq:Dunkl-small-time-smoothing}
\end{align}

It remains to treat the large-time part. Let \(K_1\Subset\Gamma_+\)
be a compact set containing \(K\). By
\eqref{eq:Dunkl-folded-heat-kernel}, the heat-kernel derivative
estimate cited above, and \eqref{doublinginequality}, for
\(x,y\in K_1\) and \(t\geq1\),
\begin{align*}
V(gx,\ell y,\sqrt t)^{-1}
\lesssim_{K_1}
t^{-N/2}.
\end{align*}
Since \(\rho\) and all its derivatives are bounded on \(K_1\), it
follows that, for every pair of multi-indices \(\alpha,\beta\),
\begin{align}
\sup_{x,y\in K_1}
\left\|
\partial_x^\alpha
\partial_y^\beta
\widetilde h_t(x,y)
\right\|
\lesssim_{\alpha,\beta,K_1}
t^{-N/2},
\qquad
t\geq1.
\label{eq:Dunkl-large-time-local-heat}
\end{align}
Since \(N>1\),
\[
\int_1^\infty
t^{-1/2}t^{-N/2}\,dt
<
\infty.
\]
Thus all derivatives of
\begin{align*}
K_{S_\infty}(x,y)
:=
\frac{\chi(x)\chi(y)}{\sqrt\pi}
\int_1^\infty
t^{-1/2}\widetilde h_t(x,y)\,dt
\end{align*}
converge locally uniformly, and therefore
\begin{align}
S_\infty
:=
\frac1{\sqrt\pi}
\int_1^\infty
t^{-1/2}
M_\chi H_tM_\chi\,dt
\in
\Psi^{-\infty}.
\label{eq:Dunkl-large-time-smoothing}
\end{align}

For completeness, the preceding integrals give the canonical bounded
localization of \(\widetilde P^{-1/2}\). Indeed, for
\(\varepsilon>0\),
\begin{align*}
(\widetilde P+\varepsilon)^{-1/2}
=
\frac1{\sqrt\pi}
\int_0^\infty
t^{-1/2}e^{-\varepsilon t}H_t\,dt
\end{align*}
by the spectral theorem. The estimates
\eqref{eq:Dunkl-local-heat-flat} and
\eqref{eq:Dunkl-large-time-local-heat} imply that
\[
M_\chi(\widetilde P+\varepsilon)^{-1/2}M_\chi
\]
converges in operator norm, as \(\varepsilon\downarrow0\), to
\begin{align*}
\frac1{\sqrt\pi}
\int_0^\infty
t^{-1/2}M_\chi H_tM_\chi\,dt.
\end{align*}
On the natural domain of \(\widetilde P^{-1/2}\), this limit agrees
with \(M_\chi\widetilde P^{-1/2}M_\chi\); hence it is its bounded
extension.

Combining
\eqref{eq:Dunkl-small-time-smoothing},
\eqref{eq:Dunkl-large-time-smoothing}, and the definition of \(B_0\),
we obtain
\begin{align*}
M_\chi\widetilde Q M_\chi
&=
M_\chi B_0M_\chi
+
S_0
+
S_\infty.
\end{align*}
By \eqref{eq:Dunkl-B0-symbol},
\(M_\chi B_0M_\chi\) is a classical matrix-valued
pseudodifferential operator of order \(-1\), while
\(S_0+S_\infty\) is smoothing. Therefore
\[
M_\chi\widetilde Q M_\chi
\in
\Psi_{\mathrm{cl}}^{-1},
\]
and, by \eqref{eq:Dunkl-B0-local-symbol},
\begin{align*}
\sigma_{-1}
\bigl(
M_\chi\widetilde Q M_\chi
\bigr)(x,\xi)
=
\chi(x)^2|\xi|^{-1}I_{\mathbb C^G}.
\end{align*}
If \(\chi=1\) in a neighborhood of \(x\), this reduces to
\[
\sigma_{-1}
\bigl(
M_\chi\widetilde Q M_\chi
\bigr)(x,\xi)
=
|\xi|^{-1}I_{\mathbb C^G},
\]
which proves \eqref{eq:Dunkl-folded-Q-principal-symbol}.
\end{proof}

\begin{lemma}\label{lem:Dunkl-commutator-principal-symbol}
Let \(j\in\{1,\ldots,N\}\) and \(f\in C_c^\infty(\mathbb R^N)\), and set
\begin{align*}
A_{f,j}
:=
VU[R_{{\rm Dunkl},j},M_f]U^{-1}V^{-1}.
\end{align*}
Then, for every \(\chi\in C_c^\infty(\Gamma_+)\),
\[
M_\chi A_{f,j}M_\chi
\]
is a compactly supported classical matrix valued pseudodifferential
operator of order \(-1\), with principal symbol
\begin{align}
\sigma_{-1}
\big(
M_\chi A_{f,j}M_\chi
\big)(x,\xi)
=
\frac{\chi(x)^2}{\|\xi\|}
\mathfrak c_{f,j}
\bigg(
x,\frac{\xi}{\|\xi\|}
\bigg),
\qquad
x\in\Gamma_+,\quad \xi\neq0.
\label{eq:Dunkl-commutator-principal-symbol}
\end{align}
\end{lemma}

\begin{proof}
Set
\begin{align*}
F(x)
:=
\operatorname{diag}_{g\in G}\big(f_g(x)\big).
\end{align*}
Since \(VU M_fU^{-1}V^{-1}=M_F\) and
\(R_{{\rm Dunkl},j}=T_{e_j}(-\Delta_{{\rm Dunkl}})^{-1/2}\),
\begin{align}
A_{f,j}
=
[\widetilde T_j,M_F]\widetilde Q
+
\widetilde T_j[\widetilde Q,M_F].
\label{eq:Dunkl-commutator-folded-decomposition}
\end{align}
By \eqref{eq:Dunkl-folded-T},
\begin{align}
[\widetilde T_j,M_F]
&=
\operatorname{diag}_{g\in G}
\big(
\langle d_{g,j},\nabla f_g\rangle
\big)
\nonumber\\
&\quad+
\sum_{\alpha\in R_+}
k(\alpha)
\operatorname{diag}_{g\in G}
\bigg(
\frac{\langle\alpha,d_{g,j}\rangle}
{\langle\alpha,x\rangle}
\big(
f_g-f_{g\sigma_\alpha}
\big)
\bigg)S_\alpha .
\label{eq:Dunkl-TF-folded-commutator}
\end{align}

Fix \(\chi\in C_c^\infty(\Gamma_+)\). Choose
\(\eta,\psi\in C_c^\infty(\Gamma_+)\) such that
\(\eta=1\) on a neighborhood of \(\operatorname{supp}\chi\) and
\(\psi=1\) on a neighborhood of \(\operatorname{supp}\eta\), and set
\begin{align*}
Q_\psi
:=
M_\psi\widetilde Q M_\psi.
\end{align*}
By Lemma \ref{lem:Dunkl-chamber-folding-symbol},
\(Q_\psi\) is a classical matrix valued pseudodifferential operator
of order \(-1\), and on a neighborhood of
\(\operatorname{supp}\eta\),
\begin{align}
q_{-1}(x,\xi)
=
\|\xi\|^{-1}I_{\mathbb C^G}.
\label{eq:Dunkl-Q-qminusone}
\end{align}

We use the Kohn--Nirenberg convention
\begin{align*}
\operatorname{Op}(a)u(x)
=
(2\pi)^{-N}
\int_{\mathbb R^{2N}}
e^{i(x-y)\cdot\xi}
a(x,\xi)u(y)\,dy\,d\xi.
\end{align*}
If \(a\) and \(b\) are classical symbols, the composition formula
\cite[Theorem~18.1.8, p.~71]{MR0781536} gives
\begin{align}
a\# b
\sim
\sum_{\beta\in\mathbb N_0^N}
\frac{1}{\beta!}
\partial_\xi^\beta a\,
D_x^\beta b,
\qquad
D_x:=\frac{1}{i}\partial_x.
\label{eq:KN-composition}
\end{align}

Write, on a neighborhood of \(\operatorname{supp}\eta\),
\begin{align*}
\sigma(Q_\psi)(x,\xi)
\sim
q_{-1}(x,\xi)+q_{-2}(x,\xi)+\cdots.
\end{align*}
By \eqref{eq:Dunkl-folded-Laplacian},
\begin{align*}
\sigma(\widetilde P)(x,\xi)
\sim
p_2(x,\xi)+p_1(x,\xi)+p_0(x,\xi),
\end{align*}
where
\begin{align*}
p_2(x,\xi)
=
\|\xi\|^2I_{\mathbb C^G},
\qquad
p_1(x,\xi)=0.
\end{align*}
Since \(\widetilde Q=\widetilde P^{-1/2}\),
\(\widetilde Q\widetilde P\widetilde Q=I\). Applying
\eqref{eq:KN-composition} on the region where \(\psi=1\) and comparing
the homogeneous terms of degree \(-1\), the terms containing
\(x\)-derivatives of \(p_2\) and \(q_{-1}\) vanish, and hence
\begin{align*}
q_{-2}p_2q_{-1}
+
q_{-1}p_2q_{-2}
=
0.
\end{align*}
Together with \eqref{eq:Dunkl-Q-qminusone}, this gives
\begin{align}
q_{-2}(x,\xi)=0
\label{eq:Dunkl-Q-qminustwo}
\end{align}
on a neighborhood of \(\operatorname{supp}\eta\).

Applying \eqref{eq:KN-composition} to the commutator
\([Q_\psi,M_F]\), and using
\eqref{eq:Dunkl-Q-qminusone}--\eqref{eq:Dunkl-Q-qminustwo}, we obtain,
on the same neighborhood,
\begin{align}
\sigma_{-2}
\big(
[Q_\psi,M_F]
\big)(x,\xi)
&=
\frac{1}{i}
\sum_{\ell=1}^N
\partial_{\xi_\ell}
\big(
\|\xi\|^{-1}
\big)
\partial_{x_\ell}F(x)
\nonumber\\
&=
i\|\xi\|^{-3}
\operatorname{diag}_{g\in G}
\big(
\langle\xi,\nabla f_g(x)\rangle
\big).
\label{eq:Dunkl-QF-principal-symbol}
\end{align}

Since \(\psi=1\) on a neighborhood of
\(\operatorname{supp}\eta\) and
\(\eta=1\) on a neighborhood of
\(\operatorname{supp}\chi\),
\begin{align*}
M_\chi[\widetilde T_j,M_F]\widetilde Q M_\chi
&=
M_\chi[\widetilde T_j,M_F]Q_\psi M_\chi,\\
M_\chi\widetilde T_j[\widetilde Q,M_F]M_\chi
&=
M_\chi\widetilde T_jM_\eta
[Q_\psi,M_F]M_\chi.
\end{align*}
Both operators on the right are classical matrix valued
pseudodifferential operators of order \(-1\).

By \eqref{eq:Dunkl-TF-folded-commutator},
\eqref{eq:Dunkl-Q-qminusone}, and the order-zero part of the
composition formula,
\begin{align}
&
\sigma_{-1}
\big(
M_\chi[\widetilde T_j,M_F]
\widetilde Q M_\chi
\big)(x,\xi)
\nonumber\\
&\quad=
\frac{\chi(x)^2}{\|\xi\|}
\operatorname{diag}_{g\in G}
\big(
\langle d_{g,j},\nabla f_g(x)\rangle
\big)
\nonumber\\
&\qquad+
\frac{\chi(x)^2}{\|\xi\|}
\sum_{\alpha\in R_+}
k(\alpha)
\operatorname{diag}_{g\in G}
\bigg(
\frac{\langle\alpha,d_{g,j}\rangle}
{\langle\alpha,x\rangle}
\big(
f_g(x)-f_{g\sigma_\alpha}(x)
\big)
\bigg)S_\alpha.
\label{eq:Dunkl-first-principal-contribution}
\end{align}

Under the above Kohn--Nirenberg convention,
\eqref{eq:Dunkl-folded-T} gives
\begin{align*}
\sigma_1(\widetilde T_j)(x,\xi)
=
i\operatorname{diag}_{g\in G}
\big(
\langle d_{g,j},\xi\rangle
\big).
\end{align*}
Therefore, by \eqref{eq:Dunkl-QF-principal-symbol},
\begin{align}
&
\sigma_{-1}
\big(
M_\chi\widetilde T_j
[\widetilde Q,M_F]M_\chi
\big)(x,\xi)
\nonumber\\
&\quad=
-\chi(x)^2\|\xi\|^{-3}
\operatorname{diag}_{g\in G}
\big(
\langle d_{g,j},\xi\rangle
\langle\xi,\nabla f_g(x)\rangle
\big).
\label{eq:Dunkl-second-principal-contribution}
\end{align}

Combining
\eqref{eq:Dunkl-first-principal-contribution} and
\eqref{eq:Dunkl-second-principal-contribution}, and setting
\begin{align*}
\theta:=\frac{\xi}{\|\xi\|},
\end{align*}
we obtain
\begin{align*}
\sigma_{-1}
\big(
M_\chi A_{f,j}M_\chi
\big)(x,\xi)
&=
\frac{\chi(x)^2}{\|\xi\|}
\operatorname{diag}_{g\in G}
\bigg(
\langle d_{g,j},\nabla f_g(x)\rangle
-
\langle d_{g,j},\theta\rangle
\langle\theta,\nabla f_g(x)\rangle
\bigg)
\\
&\quad+
\frac{\chi(x)^2}{\|\xi\|}
\sum_{\alpha\in R_+}
k(\alpha)
\operatorname{diag}_{g\in G}
\bigg(
\frac{\langle\alpha,d_{g,j}\rangle}
{\langle\alpha,x\rangle}
\big(
f_g(x)-f_{g\sigma_\alpha}(x)
\big)
\bigg)S_\alpha
\\
&=
\frac{\chi(x)^2}{\|\xi\|}
\mathfrak c_{f,j}(x,\theta).
\end{align*}
This proves \eqref{eq:Dunkl-commutator-principal-symbol}.
\end{proof}



\begin{lemma}\label{lem:Dunkl-spectral-density-Sobolev}
Assume that \(N>1\), \(j\in\{1,\ldots,N\}\), and \(f\in\dot W^{1,N}(\mathbb R^N)\). Then
\begin{align*}
\Bigg(
\int_{\Gamma}
\int_{\mathbb S^{N-1}}
\operatorname{tr}_{\mathbb C^G}
\big|
\mathfrak c_{f,j}(x,\theta)
\big|^N
\,d\theta\,dx
\Bigg)^{1/N}
\simeq_{N,R,k}
\|\nabla f\|_{L^N(\mathbb R^N,dx)}.
\end{align*}
\end{lemma}
\begin{proof}
By the definition of \(\mathfrak c_{f,j}\), the unitary invariance of
the Schatten norm, and \(\|d_{g,j}\|=1\),
\begin{align*}
\|\mathfrak c_{f,j}(x,\theta)\|_{S^N}
&\lesssim_{R,k}
\Bigg(
\sum_{g\in G}
|\nabla f_g(x)|^N
\Bigg)^{1/N}\\
&\quad+
\sum_{\alpha\in R_+}
\Bigg(
\sum_{g\in G}
\bigg|
\frac{f_g(x)-f_{g\sigma_\alpha}(x)}
{\langle\alpha,x\rangle}
\bigg|^N
\Bigg)^{1/N}.
\end{align*}
The chamber decomposition used in Lemma
\ref{lem:Dunkl-chamber-folding-symbol} and the orthogonality of
\(g\in G\) imply
\begin{align*}
\sum_{g\in G}
\int_\Gamma
|\nabla f_g(x)|^N\,dx
=
\|\nabla f\|_{L^N(\mathbb R^N,dx)}^N.
\end{align*}
Moreover, since \(\|\alpha\|^2=2\),
\begin{align*}
|x-\sigma_\alpha x|
=
\sqrt{2}\,|\langle\alpha,x\rangle|,
\end{align*}
and hence
\begin{align*}
\bigg|
\frac{f_g(x)-f_{g\sigma_\alpha}(x)}
{\langle\alpha,x\rangle}
\bigg|
=
\sqrt{2}\,
\mathcal G_{g,g\sigma_\alpha}f(x)
\end{align*}
for almost every \(x\in\mathbb R^N\). By Lemma
\ref{lem:Weyl-difference-quotient},
\begin{align*}
\int_\Gamma
\bigg|
\frac{f_g(x)-f_{g\sigma_\alpha}(x)}
{\langle\alpha,x\rangle}
\bigg|^N
\,dx
\lesssim_{N,G}
\|\nabla f\|_{L^N(\mathbb R^N,dx)}^N
\end{align*}
uniformly in \(g\in G\) and \(\alpha\in R_+\). Since \(G\) and
\(R_+\) are finite,
\begin{align}
\Bigg(
\int_\Gamma
\int_{\mathbb S^{N-1}}
\operatorname{tr}_{\mathbb C^G}
\big|
\mathfrak c_{f,j}(x,\theta)
\big|^N
\,d\theta\,dx
\Bigg)^{1/N}
\lesssim_{N,R,k}
\|\nabla f\|_{L^N(\mathbb R^N,dx)}.
\label{eq:Dunkl-spectral-density-upper}
\end{align}

For the converse estimate, let \(P_g\) denote the orthogonal projection
of \(\mathbb C^G\) onto the \(g\)-th coordinate and define
\begin{align*}
\mathcal E(A)
:=
\sum_{g\in G}P_gAP_g,
\qquad
A\in\mathcal L(\mathbb C^G).
\end{align*}
Thus, \(\mathcal E(A)\) is the diagonal part of \(A\). The pinching
inequality for unitarily invariant norms
\cite[p.~97, (IV.51)--(IV.52)]{BhatiaMatrixAnalysis} gives
\begin{align}
\|\mathcal E(A)\|_{S^N}
\leq
\|A\|_{S^N}.
\label{eq:Schatten-diagonal-contraction}
\end{align}

For every \(\alpha\in R_+\), \(S_\alpha\) has zero diagonal, since
\(g\sigma_\alpha\neq g\) for all \(g\in G\). Hence, by the definition
of \(\mathfrak c_{f,j}\),
\begin{align*}
\mathcal E\big(\mathfrak c_{f,j}(x,\theta)\big)
=
\operatorname{diag}_{g\in G}
\bigg(
\langle d_{g,j},\nabla f_g(x)\rangle
-
\langle d_{g,j},\theta\rangle
\langle\theta,\nabla f_g(x)\rangle
\bigg).
\end{align*}
Therefore, by \eqref{eq:Schatten-diagonal-contraction},
\begin{align}
\|\mathfrak c_{f,j}(x,\theta)\|_{S^N}^N
\geq
\sum_{g\in G}
\big|
\langle d_{g,j},\nabla f_g(x)\rangle
-
\langle d_{g,j},\theta\rangle
\langle\theta,\nabla f_g(x)\rangle
\big|^N.
\label{eq:Dunkl-spectral-density-diagonal-lower}
\end{align}
For each \(g\in G\), set
\begin{align*}
y:=gx,
\qquad
\eta:=g\theta.
\end{align*}
Since
\begin{align*}
\nabla f_g(x)
=
g^{-1}\nabla f(y),
\qquad
d_{g,j}=g^{-1}e_j,
\end{align*}
and \(g\) is orthogonal,
\begin{align*}
\langle d_{g,j},\nabla f_g(x)\rangle
&=
\partial_j f(y),\\
\langle d_{g,j},\theta\rangle
&=
\eta_j,\\
\langle\theta,\nabla f_g(x)\rangle
&=
\langle\eta,\nabla f(y)\rangle.
\end{align*}
Using the chamber decomposition and the orthogonal invariance of
surface measure, \eqref{eq:Dunkl-spectral-density-diagonal-lower}
therefore gives
\begin{align}
&\int_\Gamma
\int_{\mathbb S^{N-1}}
\|\mathfrak c_{f,j}(x,\theta)\|_{S^N}^N
\,d\theta\,dx
\nonumber\\
&\qquad\geq
\int_{\mathbb R^N}
\int_{\mathbb S^{N-1}}
\big|
\partial_jf(y)
-
\eta_j\langle\eta,\nabla f(y)\rangle
\big|^N
\,d\eta\,dy.
\label{eq:Dunkl-spectral-density-lower-reduction}
\end{align}

For \(v\in\mathbb C^N\), set
\begin{align*}
\Phi_j(v)
:=
\Bigg(
\int_{\mathbb S^{N-1}}
\big|
v_j-\eta_j\langle\eta,v\rangle
\big|^N
\,d\eta
\Bigg)^{1/N}.
\end{align*}
Then
\begin{align}
\Phi_j(v)
\simeq_N
\|v\|.
\label{eq:spherical-projected-gradient-norm}
\end{align}
Indeed, \(\Phi_j\) is a seminorm on \(\mathbb C^N\). If
\(\Phi_j(v)=0\), continuity gives
\begin{align*}
v_j-\eta_j\langle\eta,v\rangle=0,
\qquad
\eta\in\mathbb S^{N-1}.
\end{align*}
Since \(N>1\), there exists \(\eta\in\mathbb S^{N-1}\) with
\(\eta_j=0\), and therefore \(v_j=0\). It follows that
\begin{align*}
\langle\eta,v\rangle=0
\end{align*}
whenever \(\eta_j\neq0\). Since such \(\eta\) form a dense subset of
\(\mathbb S^{N-1}\), continuity yields
\(\langle\eta,v\rangle=0\) for every
\(\eta\in\mathbb S^{N-1}\), and hence \(v=0\). Thus \(\Phi_j\) is a
norm, and \eqref{eq:spherical-projected-gradient-norm} follows from
the equivalence of norms on the finite dimensional space
\(\mathbb C^N\).

Applying \eqref{eq:spherical-projected-gradient-norm} with
\(v=\nabla f(y)\) in
\eqref{eq:Dunkl-spectral-density-lower-reduction}, we obtain
\begin{align*}
\|\nabla f\|_{L^N(\mathbb R^N,dx)}
\lesssim_N
\Bigg(
\int_\Gamma
\int_{\mathbb S^{N-1}}
\operatorname{tr}_{\mathbb C^G}
\big|
\mathfrak c_{f,j}(x,\theta)
\big|^N
\,d\theta\,dx
\Bigg)^{1/N}.
\end{align*}
Combining this estimate with
\eqref{eq:Dunkl-spectral-density-upper} proves the assertion.
\end{proof}
\begin{lemma}\label{lem:Dunkl-smooth-Weyl-asymptotic}
Assume that \(N>1\), \(j\in\{1,\ldots,N\}\), and \(f\in C_c^\infty(\mathbb R^N)\). Then
\begin{align}
\lim_{t\to\infty}
t^{1/N}
\mu
\big(
t,[R_{{\rm Dunkl},j},M_f]
\big)
=
\Bigg(
\frac{1}{N(2\pi)^N}
\int_\Gamma
\int_{\mathbb S^{N-1}}
\operatorname{tr}_{\mathbb C^G}
\big|
\mathfrak c_{f,j}(x,\theta)
\big|^N
\,d\theta\,dx
\Bigg)^{1/N}.
\label{eq:Dunkl-smooth-Weyl-asymptotic}
\end{align}
\end{lemma}

\begin{proof}
Define
\begin{align*}
A_f
:=
VU[R_{{\rm Dunkl},j},M_f]U^{-1}V^{-1}.
\end{align*}
By Lemma \ref{lem:Dunkl-automatic-wall-localization}
\eqref{eq:Dunkl-automatic-wall-localization}, there exist
\(\chi_m\in C_c^\infty(\Gamma_+)\), \(0\leq\chi_m\leq1\), such that
\(\chi_m(x)\to1\) for every \(x\in\Gamma_+\) and
\begin{align}
\big\|
A_f-M_{\chi_m}A_fM_{\chi_m}
\big\|_{S^{N,\infty}}
\longrightarrow0.
\label{eq:Dunkl-smooth-localization-convergence}
\end{align}
Set
\[
A_{f,m}:=M_{\chi_m}A_fM_{\chi_m},
\qquad
K_m:=\operatorname{supp}\chi_m.
\]
By Lemma \ref{lem:Dunkl-commutator-principal-symbol}
\eqref{eq:Dunkl-commutator-principal-symbol},
\(A_{f,m}\) is a compactly supported classical matrix valued pseudodifferential operator of order \(-1\) on \(\Gamma_+\), with
\begin{align}
\sigma_{-1}(A_{f,m})(x,\xi)
=
\frac{\chi_m(x)^2}{\|\xi\|}
\mathfrak c_{f,j}
\bigg(
x,\frac{\xi}{\|\xi\|}
\bigg),
\qquad
x\in\Gamma_+,\quad \xi\neq0.
\label{eq:Dunkl-smooth-localized-symbol}
\end{align}

Choose a bounded open set \(U_m\) with
\(K_m\Subset U_m\Subset\Gamma_+\), and choose \(L_m>0\) such that
\[
L_m>\sup_{x,y\in U_m}\|x-y\|_\infty.
\]
For
\[
\mathbb T_m^N:=\mathbb R^N/L_m\mathbb Z^N,
\]
the quotient map restricts to a diffeomorphism
\[
\kappa_m:U_m\longrightarrow\Omega_m:=\kappa_m(U_m),
\]
which is a local isometry. Since the Schwartz kernel of \(A_{f,m}\) is supported in \(K_m\times K_m\), with respect to
\[
L^2(\Gamma;\mathbb C^G)
=
L^2(U_m;\mathbb C^G)
\oplus
L^2(\Gamma\setminus U_m;\mathbb C^G)
\]
we have
\[
A_{f,m}=B_{f,m}\oplus0
\]
for some \(B_{f,m}\) on \(L^2(U_m;\mathbb C^G)\).

Let
\[
W_m:L^2(U_m;\mathbb C^G)\longrightarrow
L^2(\Omega_m;\mathbb C^G),
\qquad
W_mu:=u\circ\kappa_m^{-1}.
\]
Then \(W_m\) is unitary. By the invariance of classical pseudodifferential operators under smooth coordinate changes
\cite[Theorem~18.1.17]{MR0781536},
\[
T_{f,m}:=W_mB_{f,m}W_m^{-1}
\in
\Psi_{\mathrm{cl}}^{-1}
\big(
\Omega_m;\mathbb C^G
\big).
\]
Its Schwartz kernel is supported in
\(\kappa_m(K_m)\times\kappa_m(K_m)\Subset\Omega_m\times\Omega_m\).
Choose \(\vartheta_m\in C_c^\infty(\Omega_m)\) equal to \(1\) on a neighborhood of \(\kappa_m(K_m)\). Then
\[
T_{f,m}=M_{\vartheta_m}T_{f,m}M_{\vartheta_m}.
\]
Hence the zero extension of the Schwartz kernel of \(T_{f,m}\) from
\(\Omega_m\times\Omega_m\) to
\(\mathbb T_m^N\times\mathbb T_m^N\) is a classical pseudodifferential kernel of order \(-1\): locally in \(\Omega_m\times\Omega_m\) this follows from \(T_{f,m}\in\Psi_{\mathrm{cl}}^{-1}\), while near the boundary of \(\Omega_m\times\Omega_m\) the kernel vanishes. Denote the corresponding operator on \(\mathbb T_m^N\) by \(\widehat A_{f,m}\). Thus,
\begin{align}
\widehat A_{f,m}
\in
\Psi_{\mathrm{cl}}^{-1}
\big(
\mathbb T_m^N;\mathbb C^G
\big).
\label{eq:Dunkl-smooth-torus-Psi}
\end{align}

Since \(\kappa_m\) is a local isometry, by
\eqref{eq:Dunkl-smooth-localized-symbol} and the transformation law for principal symbols,
\begin{align}
\sigma_{-1}(\widehat A_{f,m})
\big(
\kappa_m(x),\xi
\big)
=
\frac{\chi_m(x)^2}{\|\xi\|}
\mathfrak c_{f,j}
\bigg(
x,\frac{\xi}{\|\xi\|}
\bigg),
\qquad
x\in U_m,\quad \xi\neq0,
\label{eq:Dunkl-smooth-torus-symbol}
\end{align}
and the principal symbol vanishes outside \(\Omega_m\).

With respect to
\[
L^2(\mathbb T_m^N;\mathbb C^G)
=
L^2(\Omega_m;\mathbb C^G)
\oplus
L^2(\mathbb T_m^N\setminus\Omega_m;\mathbb C^G),
\]
we have
\[
\widehat A_{f,m}=T_{f,m}\oplus0.
\]
Since \(T_{f,m}\) and \(B_{f,m}\) are unitarily equivalent,
\begin{align}
\mu(t,\widehat A_{f,m})
=
\mu(t,A_{f,m}),
\qquad
t>0.
\label{eq:Dunkl-smooth-singular-values}
\end{align}

By \eqref{eq:Dunkl-smooth-torus-Psi} and the Weyl law
\cite[Theorem~6.1]{MR4592885},
\begin{align*}
\lim_{t\to\infty}
t^{1/N}\mu(t,\widehat A_{f,m})
=
\Bigg(
\frac{1}{N(2\pi)^N}
\int_{S^*\mathbb T_m^N}
\operatorname{tr}_{\mathbb C^G}
\big|
\sigma_{-1}(\widehat A_{f,m})(z,\theta)
\big|^N
\,dz\,d\theta
\Bigg)^{1/N}.
\end{align*}
Since the metric on \(\mathbb T_m^N\) is flat, combining
\eqref{eq:Dunkl-smooth-torus-symbol} with
\eqref{eq:Dunkl-smooth-singular-values}, we obtain
\begin{align}
\lim_{t\to\infty}
t^{1/N}\mu(t,A_{f,m})
=
\Bigg(
\frac{1}{N(2\pi)^N}
\int_\Gamma
\int_{\mathbb S^{N-1}}
\chi_m(x)^{2N}
\operatorname{tr}_{\mathbb C^G}
\big|
\mathfrak c_{f,j}(x,\theta)
\big|^N
\,d\theta\,dx
\Bigg)^{1/N},
\label{eq:Dunkl-smooth-localized-Weyl}
\end{align}
where \(\chi_m\) is identified with its zero extension to \(\Gamma\).

By Lemma \ref{lem:Dunkl-spectral-density-Sobolev},
\[
(x,\theta)
\longmapsto
\operatorname{tr}_{\mathbb C^G}
\big|
\mathfrak c_{f,j}(x,\theta)
\big|^N
\]
belongs to \(L^1(\Gamma\times\mathbb S^{N-1})\). Since
\(0\leq\chi_m\leq1\), \(\chi_m\to1\) on \(\Gamma_+\), and
\(\Gamma\setminus\Gamma_+\) is contained in finitely many hyperplanes, the dominated convergence theorem implies
\begin{align}
&\lim_{m\to\infty}
\Bigg(
\frac{1}{N(2\pi)^N}
\int_\Gamma
\int_{\mathbb S^{N-1}}
\chi_m(x)^{2N}
\operatorname{tr}_{\mathbb C^G}
\big|
\mathfrak c_{f,j}(x,\theta)
\big|^N
\,d\theta\,dx
\Bigg)^{1/N}
\nonumber\\
&\qquad=
\Bigg(
\frac{1}{N(2\pi)^N}
\int_\Gamma
\int_{\mathbb S^{N-1}}
\operatorname{tr}_{\mathbb C^G}
\big|
\mathfrak c_{f,j}(x,\theta)
\big|^N
\,d\theta\,dx
\Bigg)^{1/N}.
\label{eq:Dunkl-smooth-symbol-limit}
\end{align}

By \eqref{eq:Dunkl-smooth-localization-convergence},
\(A_{f,m}\to A_f\) in \(S^{N,\infty}\). Combining
\eqref{eq:Dunkl-smooth-localized-Weyl},
\eqref{eq:Dunkl-smooth-symbol-limit}, and Lemma
\ref{lem:weak-Schatten-asymptotic-stability}, we obtain
\begin{align*}
\lim_{t\to\infty}
t^{1/N}\mu(t,A_f)
=
\Bigg(
\frac{1}{N(2\pi)^N}
\int_\Gamma
\int_{\mathbb S^{N-1}}
\operatorname{tr}_{\mathbb C^G}
\big|
\mathfrak c_{f,j}(x,\theta)
\big|^N
\,d\theta\,dx
\Bigg)^{1/N}.
\end{align*}
Since \(U\) and \(V\) are unitary,
\[
\mu(t,A_f)
=
\mu
\big(
t,[R_{{\rm Dunkl},j},M_f]
\big),
\qquad
t>0.
\]
This completes the proof of Lemma \ref{lem:Dunkl-smooth-Weyl-asymptotic}.
\end{proof}

\begin{lemma}\label{lem:weak-Schatten-asymptotic-stability}
Let \(0<p<\infty\), let \(\mathcal H\) be a Hilbert space, and let
\(A,A_m\in S^{p,\infty}(\mathcal H)\), \(m\geq1\), satisfy
\begin{align*}
\|A_m-A\|_{S^{p,\infty}}
\longrightarrow0.
\end{align*}
Suppose that, for every \(m\geq1\), the limit
\begin{align*}
c_m
:=
\lim_{t\to\infty}
t^{1/p}\mu(t,A_m)
\end{align*}
exists. Then both limits
\begin{align*}
\lim_{m\to\infty}c_m
\qquad\text{and}\qquad
\lim_{t\to\infty}
t^{1/p}\mu(t,A)
\end{align*}
exist and
\begin{align*}
\lim_{t\to\infty}
t^{1/p}\mu(t,A)
=
\lim_{m\to\infty}c_m.
\end{align*}
\end{lemma}

\begin{proof}
The assertion is the Birman--Solomyak approximation lemma; see \cite[Section~11.6]{MR1192782}.
\end{proof}

\begin{lemma}\label{lem:Dunkl-Weyl-Sobolev-approximation}
Assume that \(N>1\), \(j\in\{1,\ldots,N\}\), and \(f\in\dot W^{1,N}(\mathbb R^N)\). Then there exists \(f_m\in C_c^\infty(\mathbb R^N)\) such that
\begin{align*}
\|\nabla(f_m-f)\|_{L^N(\mathbb R^N,dx)}
&\longrightarrow0,\\
\big\|
[R_{{\rm Dunkl},j},M_{f_m}]
-
[R_{{\rm Dunkl},j},M_f]
\big\|_{S^{N,\infty}}
&\longrightarrow0,\\
\|\mathfrak c_{f_m,j}-\mathfrak c_{f,j}\|_
{L^N(\Gamma\times\mathbb S^{N-1};S^N(\mathbb C^G))}
&\longrightarrow0.
\end{align*}
\end{lemma}

\begin{proof}
By the density of \(C_c^\infty(\mathbb R^N)\) in \(\dot W^{1,N}(\mathbb R^N)\), choose \(f_m\in C_c^\infty(\mathbb R^N)\) such that
\begin{align*}
\|\nabla(f_m-f)\|_{L^N(\mathbb R^N,dx)}
\longrightarrow0.
\end{align*}
By Theorem \ref{lem:Dunkl-critical-weak-Schatten-upper} and Lemma \ref{lem:Euclidean-OSC-critical-rigidity},
\begin{align*}
\big\|
[R_{{\rm Dunkl},j},M_{f_m}]
-
[R_{{\rm Dunkl},j},M_f]
\big\|_{S^{N,\infty}}
&=
\big\|
[R_{{\rm Dunkl},j},M_{f_m-f}]
\big\|_{S^{N,\infty}}\\
&\lesssim
\|f_m-f\|_{OSC^{N,\infty}(\mathbb R^N)}
\lesssim
\|\nabla(f_m-f)\|_{L^N(\mathbb R^N,dx)}
\longrightarrow0.
\end{align*}
By the linearity of \(f\mapsto\mathfrak c_{f,j}\) and Lemma \ref{lem:Dunkl-spectral-density-Sobolev},
\begin{align*}
\|\mathfrak c_{f_m,j}-\mathfrak c_{f,j}\|_
{L^N(\Gamma\times\mathbb S^{N-1};S^N(\mathbb C^G))}
&=
\|\mathfrak c_{f_m-f,j}\|_
{L^N(\Gamma\times\mathbb S^{N-1};S^N(\mathbb C^G))}\\
&\lesssim_{N,R,k}
\|\nabla(f_m-f)\|_{L^N(\mathbb R^N,dx)}
\longrightarrow0.
\end{align*}
This completes the proof of Lemma \ref{lem:Dunkl-Weyl-Sobolev-approximation}.
\end{proof}

\begin{theorem}\label{thm:Dunkl-Riesz-spectral-asymptotic}
Assume that \(N>1\) and \(j\in\{1,\ldots,N\}\). Then, for every
\(f\in\dot W^{1,N}(\mathbb R^N)\),
\begin{align}
\lim_{t\to\infty}
t^{1/N}
\mu
\big(
t,[R_{{\rm Dunkl},j},M_f]
\big)
=
\Bigg(
\frac{1}{N(2\pi)^N}
\int_\Gamma
\int_{\mathbb S^{N-1}}
\operatorname{tr}_{\mathbb C^G}
\big|
\mathfrak c_{f,j}(x,\theta)
\big|^N
\,d\theta\,dx
\Bigg)^{1/N}.
\label{eq:Dunkl-Riesz-spectral-asymptotic}
\end{align}
Moreover,
\begin{align}
\Bigg(
\int_\Gamma
\int_{\mathbb S^{N-1}}
\operatorname{tr}_{\mathbb C^G}
\big|
\mathfrak c_{f,j}(x,\theta)
\big|^N
\,d\theta\,dx
\Bigg)^{1/N}
\simeq_{N,R,k}
\|\nabla f\|_{L^N(\mathbb R^N,dx)}.
\label{eq:Dunkl-Riesz-spectral-coefficient-Sobolev}
\end{align}
\end{theorem}

\begin{proof}
By Lemma \ref{lem:Dunkl-Weyl-Sobolev-approximation}, choose
\(f_m\in C_c^\infty(\mathbb R^N)\) such that
\begin{align}
\big\|
[R_{{\rm Dunkl},j},M_{f_m}]
-
[R_{{\rm Dunkl},j},M_f]
\big\|_{S^{N,\infty}}
&\longrightarrow0,
\label{eq:Dunkl-Riesz-spectral-operator-approximation}\\
\|\mathfrak c_{f_m,j}-\mathfrak c_{f,j}\|_
{L^N(\Gamma\times\mathbb S^{N-1};S^N(\mathbb C^G))}
&\longrightarrow0.
\label{eq:Dunkl-Riesz-spectral-symbol-approximation}
\end{align}
By Lemma \ref{lem:Dunkl-smooth-Weyl-asymptotic}, for every \(m\),
\begin{align*}
\lim_{t\to\infty}
t^{1/N}
\mu
\big(
t,[R_{{\rm Dunkl},j},M_{f_m}]
\big)
=
\frac{1}{(N(2\pi)^N)^{1/N}}
\|\mathfrak c_{f_m,j}\|_
{L^N(\Gamma\times\mathbb S^{N-1};S^N(\mathbb C^G))}.
\end{align*}
By \eqref{eq:Dunkl-Riesz-spectral-symbol-approximation},
\begin{align*}
\|\mathfrak c_{f_m,j}\|_
{L^N(\Gamma\times\mathbb S^{N-1};S^N(\mathbb C^G))}
\longrightarrow
\|\mathfrak c_{f,j}\|_
{L^N(\Gamma\times\mathbb S^{N-1};S^N(\mathbb C^G))}.
\end{align*}
Combining this convergence with
\eqref{eq:Dunkl-Riesz-spectral-operator-approximation}, we apply
Lemma \ref{lem:weak-Schatten-asymptotic-stability} with \(p=N\) and obtain
\eqref{eq:Dunkl-Riesz-spectral-asymptotic}. Finally,
\eqref{eq:Dunkl-Riesz-spectral-coefficient-Sobolev} follows from
Lemma \ref{lem:Dunkl-spectral-density-Sobolev}. This completes the proof of
Theorem \ref{thm:Dunkl-Riesz-spectral-asymptotic}.
\end{proof}


\appendix

\section{Zhijie lemma}
\setcounter{equation}{0}

\begin{lemma}\label{lem:Cwikel-dyadic-distribution}
Let $1<s<\infty$, and let $\mathbb D^0$ be the standard dyadic grid
in $\mathbb R^N$. Suppose that
$
F\in L^s(\mathbb R^N).
$
Then there exists a constant $C>0$ such that
\[
\left\|\left\{|Q|^{-1/s'}\int_Q |F(x)|\,dx\right\}_{Q\in\mathbb D^0}\right\|_{\ell^{s,\infty}}
\leq C
\|F\|_{L^s(\mathbb R^N)}.
\]
\end{lemma}
\begin{proof}
Let $\mathcal E\subset\mathbb D^0$ be finite and set
\[
G_{\mathcal E}(x)
:=
\sum_{Q\in\mathcal E}|Q|^{-1/s'}\chi_Q(x).
\]
Since the dyadic cubes containing any fixed $x$ are totally ordered by inclusion, for each
$x\in\bigcup_{Q\in\mathcal E}Q$ we may choose a cube $Q_x\in\mathcal E$ of minimal volume such that $x\in Q_x$. Every $Q\in\mathcal E$ containing $x$ is then a dyadic ancestor of $Q_x$. Since the volume increases by the factor $2^N$ at each generation,
\[
G_{\mathcal E}(x)
=
\sum_{\substack{Q\in\mathcal E\\ x\in Q}}|Q|^{-1/s'}
\leq
|Q_x|^{-1/s'}
\sum_{k=0}^{\infty}2^{-kN/s'}
\lesssim_{N,s}
\sup_{Q\in\mathcal E}|Q|^{-1/s'}\chi_Q(x).
\]
Consequently,
\begin{align*}
\|G_{\mathcal E}\|_{L^{s'}(\mathbb R^N)}^{s'}
\lesssim_{N,s}
\int_{\mathbb R^N}
\sup_{Q\in\mathcal E}|Q|^{-1}\chi_Q(x)\,dx\leq
\int_{\mathbb R^N}
\sum_{Q\in\mathcal E}|Q|^{-1}\chi_Q(x)\,dx
=
\#\mathcal E.
\end{align*}
and hence,
\[
\|G_{\mathcal E}\|_{L^{s'}(\mathbb R^N)}
\lesssim_{N,s}
(\#\mathcal E)^{1/s'}.
\]
By H\"older's inequality,
\begin{align*}
\sum_{Q\in\mathcal E}|Q|^{-1/s'}\int_Q|F(x)|\,dx
=
\int_{\mathbb R^N}|F(x)|G_{\mathcal E}(x)\,dx\lesssim_{N,s}
\|F\|_{L^s(\mathbb R^N)}(\#\mathcal E)^{1/s'}.
\end{align*}

Fix $t>0$ and let $\mathcal E$ be any finite subcollection of the dyadic cubes satisfying
\[
|Q|^{-1/s'}\int_Q|F(x)|\,dx>t.
\]
Then
\[
t\,\#\mathcal E
<
\sum_{Q\in\mathcal E}|Q|^{-1/s'}\int_Q|F(x)|\,dx
\lesssim_{N,s}
\|F\|_{L^s(\mathbb R^N)}(\#\mathcal E)^{1/s'},
\]
so
\[
\#\mathcal E
\lesssim_{N,s}
t^{-s}\|F\|_{L^s(\mathbb R^N)}^s.
\]
Since $\mathcal E$ is arbitrary,
\[
\#\Bigg\{
Q\in\mathbb D^0:
|Q|^{-1/s'}\int_Q|F(x)|\,dx>t
\Bigg\}
\lesssim_{N,s}
t^{-s}\|F\|_{L^s(\mathbb R^N)}^s,
\qquad t>0.
\]
The asserted weak $\ell^s$ estimate follows. This completes the proof of Lemma~\ref{lem:Cwikel-dyadic-distribution}.
\end{proof}


\bibliographystyle{plain}
\bibliography{references}

\end{document}
